% Chapter 4, Section 1 _Linear Algebra_ Jim Hefferon
%  http://joshua.smcvt.edu/linearalgebra
%  2001-Jun-12
\chapter{行列式} 
在第一章中，我们
重点讨论了方程个数与未知量个数相同的线性方程组
这一特殊情形，即形如 \( T\vec{x}=\vec{b} \)、其中 $T$ 为方阵的方程组。
我们注意到，这样的 $T$ 只有两种。
若 $T$ 对任意 $\vec{b}$，
都对应唯一解，比如
齐次方程组 $T\vec{x}=\zero$ 那样，那么
对每一个这样的 $\vec{b}$，$T$ 都对应唯一解。
我们称这样的矩阵为非奇异的。
另一类 $T$，凡是以其为
系数矩阵的线性方程组都要么无解、
要么有无穷多解，
我们称之为奇异的。

在此后的工作中，这一区分一直是一条主线。
例如，我们现在知道，一个
\( \nbyn{n} \) 矩阵 \( T \) 是非奇异的当且仅当
以下每一条都成立：
%<*EquivalentOfNonsingular>
\begin{itemize}
   \item 任意方程组 \( T\vec{x}=\vec{b} \) 都有解，
          且该解唯一；
    \item 对 $T$ 作高斯–若尔当消元得到单位矩阵；
    \item $T$ 的各行构成一个线性无关的集合；
    \item \( T \) 的各列构成一个线性无关的集合，
          即 \( \Re^n \) 的一个基；
    \item \( T \) 所表示的任意映射都是同构；
    \item 存在逆矩阵 \( T^{-1} \)。
\end{itemize}
%</EquivalentOfNonsingular>
因此，当我们看到一个方阵时，首先
要问的问题之一就是它是否
非奇异。

本章将建立一个公式，用来判定 $T$ 是否非奇异。
更准确地说，我们将对 $\nbyn{1}$~矩阵
给出一个公式，对 $\nbyn{2}$~矩阵给出一个，等等。
这些公式自然相互关联；也就是说，
我们将建立一族
公式，一个描述每种规格之公式的方案。

由于我们的讨论仅限于方阵，本章中
我们常常径直说“矩阵”，以代替“方阵”。





\section{定义}
%<*DeterminantIntro>
判定非奇异性这件事，
对 \( \nbyn{1} \) 矩阵而言是平凡的。
\begin{equation*}
  \begin{mat}
       a
  \end{mat} 
  \quad\text{是非奇异的当且仅当}\quad
  a \neq 0
\end{equation*} 
推论~Three.IV.\ref{cor:TwoByTwoInv} 给出了 $\nbyn{2}$ 的公式。 
\begin{equation*}
     \begin{mat}
           a  &b  \\
           c  &d
      \end{mat}  
   \quad\text{是非奇异的当且仅当}\quad
   ad-bc \neq 0
\end{equation*}
我们可以像前面那样得出 $\nbyn{3}$ 的公式，
尽管计算相当繁琐
% \typeout{DET1 ThreeByThreeDetForm cleveref expmt!}
% (see \cref{exer:ThreeByThreeDetForm}).
（见\nearbyexercise{exer:ThreeByThreeDetForm}）。
\begin{equation*}
     \begin{mat}
           a  &b  &c  \\
           d  &e  &f  \\
           g  &h  &i
     \end{mat}
  \quad\text{是非奇异的当且仅当}\quad
    aei+bfg+cdh-hfa-idb-gec \neq 0
\end{equation*}
心怀这些特例，我们设想一族
公式：$a$，$ad-bc$，等等。
对每个 $n$，该公式定义了一个 
\definend{行列式}\index{determinant}\index{matrix!determinant} 
函数
$\map{\det_{\nbyn{n}}}{\matspace_{\nbyn{n}}}{\Re}$ 
使得 $\nbyn{n}$ 矩阵 $T$ 是非奇异的当
且仅当 $\det_{\nbyn{n}}(T)\neq 0$。
%</DeterminantIntro>
（我们通常省略下标 \( \nbyn{n} \)，因为
\( T \) 的规格已经表明我们指的是哪个行列式函数。）






\subsectionoptional{探索}
\textit{本小节为选读内容，
给出一般定义的动机与展开。
定义在下一小节中。}

在前文中，各种情形下的矩阵都是非奇异的当且
仅当某个公式不为零。
但这三个公式并
没有显示出明显的规律。
我们或许会注意到 \(\nbyn{1}\) 的项
\( a \) 有一个字母，\(\nbyn{2}\) 的项
\(ad\) 与 \(bc\) 各有两个字母，而 \(\nbyn{3}\)
的项各有三个字母。
我们甚至能注意到，在这些项中
每一项都是从矩阵的每一行与每一列各取一个字母，例如，
在 \(cdh\) 项中一个字母
来自每一行、每一列。
\begin{equation*}
   \begin{mat}
          &    &c \\
      d           \\
          &h
   \end{mat} 
\end{equation*}
但这些观察与其说是有启发性
的，不如说是令人困惑的。
比如，我们可能想知道为什么
有些项相加，而有些项相减。

解题的一个好策略是
先探索解必须具有哪些性质，然后
再去寻找具有这些性质的东西。
于是我们先来问：我们希望行列式
公式具有哪些性质。

就目前而言，
判定一个矩阵是否奇异的主要方法
是作高斯
消元，然后检查
阶梯形矩阵的对角线上是否有零，
即沿对角线的乘积是否为~零。
于是我们可以猜想：无论找到怎样的行列式公式，证明
它正确的做法可能都涉及对矩阵应用高斯消元法，
以说明最终沿对角线的乘积为零当且仅当
我们的公式给出零。

这提示了一个计划：~我们将寻找一族行列式
公式，它们
不受行变换影响，且阶梯形
矩阵的行列式是其对角元之积。
% Under this plan, a proof that the functions determine singularity would go, 
% ``Where $T\rightarrow\cdots\rightarrow\hat{T}$ is the Gaussian
% reduction, the determinant of $T$ equals the
% determinant of $\hat{T}$ (because the determinant is unchanged by row
% operations), which is the product down the diagonal, which is
% zero if and only if the matrix is singular''.
在本小节余下的部分，我们将用 $\nbyn{2}$ 与
$\nbyn{3}$ 的公式来检验这一计划。
最后我们将不得不修改“不受行变换影响”
这一部分，但改动不大。

首先我们检查
$\nbyn{2}$ 与 $\nbyn{3}$ 的公式是否不受
倍加这一行变换的影响：~若
\begin{equation*}
   T \grstep{k\rho_i+\rho_j} \hat{T}
\end{equation*}
那么 \( \det(\hat{T})=\det(T) \) 是否成立？ 
对 $k\rho_1+\rho_2$ 操作之后的 $\nbyn{2}$ 行列式作这一检验
\begin{equation*}
   \det(
       \begin{mat}
         a     &b       \\
         ka+c  &kb+d    \\
      \end{mat}
   )
   = a(kb+d)-(ka+c)b = ad-bc
\end{equation*}
表明它确实不变，而
另一个 $\nbyn{2}$ 倍加 $k\rho_2+\rho_1$ 给出同样的结果。
同样，
$\nbyn{3}$ 的倍加 $k\rho_3+\rho_2$ 也使行列式保持不变
  \begin{align*}
    \det(
      \begin{mat}
         a    &b    &c    \\
         kg+d &kh+e &ki+f \\
         g    &h    &i
      \end{mat}
    )
    &=\begin{array}[t]{@{}l@{}}
         a(kh+e)i+b(ki+f)g+c(kg+d)h \\
         \ \hbox{}-h(ki+f)a-i(kg+d)b-g(kh+e)c  
    \end{array}                                 \\
     &=aei + bfg + cdh - hfa - idb - gec
  \end{align*}
其他的 $\nbyn{3}$ 行倍加操作也是如此。

所以这个计划看起来是有希望的。
当然，也许如果我们算出
$\nbyn{4}$ 行列式公式并加以检验，可能会发现
它会受行倍加的影响。
这只是一次探索，我们还没有掌握全部事实。
尽管如此，到目前为止一切顺利。

接下来，对于行对换的情形，我们比较 \( \det(\hat{T}) \) 与
\( \det(T) \)。
% \begin{equation*}
%    T \grstep{ {\rho}_i \leftrightarrow {\rho}_j } \hat{T}
% \end{equation*}
这里我们遇到了一个障碍：
\(\nbyn{2}\) 的行对换 $\rho_1\leftrightarrow\rho_2$
并不给出 \( ad-bc \)。
  \begin{equation*}
     \det(
       \begin{mat}
         c  &d \\
         a  &b
       \end{mat}
     )
     = bc - ad
  \end{equation*}
而 \(\nbyn{3}\) 矩阵内的这次 $\rho_1\leftrightarrow\rho_3$ 对换
\begin{equation*}
   \det(
     \begin{mat}
        g  &h  &i \\
        d  &e  &f \\
        a  &b  &c
     \end{mat}
   )
   = gec + hfa + idb - bfg - cdh - aei
\end{equation*}
同样不给出与对换前相同的行列式，因为又一次
出现了符号改变。
再试 \(\nbyn{3}\) 的另一次对换 $\rho_1\leftrightarrow\rho_2$ 
\begin{equation*}
   \det(
     \begin{mat}
        d  &e  &f \\
        a  &b  &c \\
        g  &h  &i
     \end{mat}
   )
   = dbi + ecg + fah - hcd - iae - gbf
\end{equation*}
也给出符号的改变。

所以在这一实验中，行对换
看来会改变行列式的符号。
这并没有完全破坏我们的计划。
我们希望通过只考虑
公式是否给出零来判定非奇异性，而不考虑其符号。
因此，与其期望行列式公式
完全不受行变换影响，不如修改计划：在行对换时
它们改变符号。

显然，
我们最后要比较 \( \det(\hat{T}) \) 与 \( \det(T) \) 
针对的操作
% \begin{equation*}
%    T \grstep{ k{\rho}_i } \hat{T}
% \end{equation*}
是把一行乘以一个标量。 
下面这个计算
\begin{equation*}
  \det(
    \begin{mat}
      a   &b   \\
      kc  &kd
    \end{mat}
  )
  = a(kd) - (kc)b
  =k\cdot (ad-bc)
\end{equation*}
最后整个行列式乘上了~$k$，
另一个 $\nbyn{2}$ 情形也得到同样的结果。
这个 \(\nbyn{3}\) 情形以同样的方式收尾
\begin{align*}
  \det(
  \begin{mat}
    a    &b    &c   \\
    d    &e    &f   \\
    kg   &kh   &ki
  \end{mat}
  )
   &= \begin{array}[t]{@{}l@{}}
         ae(ki) + bf(kg) + cd(kh)                \\
         \>- (kh)fa - (ki)db - (kg)ec  
      \end{array}                                      \\
   &= k\cdot(aei + bfg + cdh - hfa - idb - gec)
\end{align*}
另外两个 $\nbyn{3}$ 情形也是如此。
这些使我们猜想：把一行乘以~$k$
会把行列式乘以~$k$。
与前面一样，这修改了我们的计划，但没有破坏它。
我们只要求
行列式公式的为零性保持不变，而不关注
它的符号或大小。

所以在这次探索中，我们的计划
经过一些非实质性的修改，现在变为：~我们将寻找
$\nbyn{n}$ 行列式
函数，它们在行倍加操作下
保持不变，在行对换时
改变符号，在某行被倍乘时按比例缩放，
并且阶梯形矩阵的行列式
是沿对角线的乘积。
在下面两小节中我们将看到，对每个~$n$，
这样的函数有且仅有一个。

最后，为下一小节作准备，请注意提出标量因子是一种
逐行进行的操作：
在这里
\begin{equation*}
    \det(
      \begin{mat}[r]
        3  &3  &9  \\
        2  &1  &1  \\
        5  &11 &-5
     \end{mat}
    )
    =3 \cdot \det(
               \begin{mat}[r]
                 1  &1  &3  \\
                 2  &1  &1  \\
                 5  &11 &-5
               \end{mat}
             )                                  
\end{equation*}
$3$ 仅从最上面一行提出，其余各行保持不变。
因此，在行列式的定义中，我们将
把它写成各行的函数
\( \det (\vec{\rho}_1,\vec{\rho}_2,\dots\vec{\rho}_n) \)，而不是
\( \det(T) \)，也不是各元素
\( \det(t_{1,1},\dots,t_{n,n}) \) 的函数。

\begin{exercises}
  \recommended \item 
     求下列各矩阵的行列式。
     \begin{exparts*}
       \partsitem \(
            \begin{mat}[r]
                3    &1   \\
               -1    &1
            \end{mat}    \)
       \partsitem \(
            \begin{mat}[r]
                2    &0   &1  \\
                3    &1   &1 \\
               -1    &0   &1
            \end{mat}    \)
       \partsitem \(
            \begin{mat}[r]
                4    &0   &1  \\
                0    &0   &1 \\
                1    &3   &-1
            \end{mat}    \)
     \end{exparts*}
     \begin{answer}
       \begin{exparts*}
         \partsitem \( 4 \)
         \partsitem \( 3 \)
         \partsitem \( -12 \)
       \end{exparts*}  
     \end{answer}
  \item 
     求下列各矩阵的行列式。
     \begin{exparts*}
        \partsitem \( \begin{mat}[r]
                    2  &0  \\
                   -1  &3
                 \end{mat} \)
        \partsitem \( \begin{mat}[r]
                    2  &1  &1  \\
                    0  &5  &-2 \\
                    1  &-3 &4
                 \end{mat} \)
        \partsitem \( \begin{mat}[r]
                    2  &3  &4  \\
                    5  &6  &7  \\
                    8  &9  &1
                 \end{mat} \)
     \end{exparts*}
    \begin{answer}
      \begin{exparts*}
        \partsitem \( 6 \)
        \partsitem \( 21 \)
        \partsitem \( 27 \)
      \end{exparts*}  
    \end{answer}
  \recommended \item  
    验证上三角的
    $\nbyn{3}$ 矩阵的行列式是对角元之积。
    \begin{equation*}
       \det(
       \begin{mat}
           a    &b   &c    \\
           0    &e   &f    \\
           0    &0   &i
       \end{mat}
       )
       =aei
    \end{equation*}
    下三角矩阵也是如此吗？
    \begin{answer}
      对于第一个，应用本节中的公式，注意任何
      含有 \( d \)、\( g \) 或 \( h \) 的项都是零，再化简即可。
      下三角矩阵也是如此。  
    \end{answer}
  \recommended \item 
     用行列式判断下列各矩阵是奇异的还是
     非奇异的。
     \begin{exparts*}
       \partsitem \(
            \begin{mat}[r]
                2    &1   \\
                3    &1
            \end{mat}    \)
       \partsitem \(
            \begin{mat}[r]
                0    &1   \\
                1    &-1
            \end{mat}    \)
       \partsitem \(
            \begin{mat}[r]
                4    &2   \\
                2    &1
            \end{mat}    \)
     \end{exparts*}
     \begin{answer}
       \begin{exparts}
         \partsitem 非奇异的，行列式为 \( -1 \)。
         \partsitem 非奇异的，行列式为 \( -1 \)。
         \partsitem 奇异的，行列式为 \( 0 \)。
       \end{exparts}  
     \end{answer}
  \item 
     奇异的还是非奇异的？
     用行列式来判断。
     \begin{exparts*}
       \partsitem \(
            \begin{mat}[r]
                2    &1   &1  \\
                3    &2   &2 \\
                0    &1   &4
            \end{mat}    \)
       \partsitem \(
            \begin{mat}[r]
                1    &0   &1  \\
                2    &1   &1 \\
                4    &1   &3
            \end{mat}    \)
       \partsitem \(
            \begin{mat}[r]
                2    &1   &0  \\
                3    &-2  &0 \\
                1    &0   &0
            \end{mat}    \)
     \end{exparts*}
     \begin{answer}
      \begin{exparts}
         \partsitem 非奇异的，行列式为 \( 3 \)。
         \partsitem 奇异的，行列式为 \( 0 \)。
         \partsitem 奇异的，行列式为 \( 0 \)。
       \end{exparts}  
     \end{answer}
  \recommended \item
    下面每一对矩阵仅相差一次行变换。
    请利用这次行变换
    来比较 \( \det(A) \) 与 \( \det(B) \)。
     \begin{exparts}
        \partsitem \( A=\begin{mat}[r]
                      1  &2  \\
                      2  &3
                   \end{mat} \), 
               \(  B=\begin{mat}[r]
                      1  &2  \\
                      0  &-1
                   \end{mat}  \)
        \partsitem \( A=\begin{mat}[r]
                      3  &1  &0  \\
                      0  &0  &1  \\
                      0  &1  &2
                   \end{mat} \),
              \(   B=\begin{mat}[r]
                      3  &1  &0  \\
                      0  &1  &2  \\
                      0  &0  &1
                   \end{mat}  \)
        \partsitem \( A=\begin{mat}[r]
                      1  &-1 &3  \\
                      2  &2  &-6 \\
                      1  &0  &4
                   \end{mat} \),
              \(   B=\begin{mat}[r]
                      1  &-1 &3  \\
                      1  &1  &-3 \\
                      1  &0  &4
                   \end{mat}  \)
     \end{exparts}
     \begin{answer}
       \begin{exparts}
         \partsitem \( \det(B)=\det(A) \)，通过 \( -2\rho_1+\rho_2 \)
         \partsitem \( \det(B)=-\det(A) \)，通过 
             \( \rho_2\leftrightarrow\rho_3 \)
         \partsitem \( \det(B)=(1/2)\cdot \det(A) \)，通过 \( (1/2)\rho_2 \)
       \end{exparts}   
     \end{answer}
  \recommended \item 按如下方案求这个 $\nbyn{4}$ 矩阵的行列式：~先做高斯消元，
    并留意行列式在倍加时保持不变，在
    对换两行时变号，在某行倍乘时相应倍乘，
    且阶梯形矩阵的行列式等于
    其对角元之积。
    \begin{equation*}
      \begin{mat}
        1 &2 &0  &2 \\
        2 &4 &1  &0 \\
        0 &0 &-1 &3 \\
        3 &-1 &1 &4
      \end{mat}
    \end{equation*}
    \begin{answer}
      高斯消元法正是这样做的。
      % [ 0 -7  1 -2]
      %  swap row 2 with row 4
      % [ 1  2  0  2]
      % [ 0 -7  1 -2]
      % [ 0  0 -1  3]
      % [ 0  0  1 -4]
      %  take 1 times row 3 plus row 4
      % [ 1  2  0  2]
      % [ 0 -7  1 -2]
      % [ 0  0 -1  3]
      % [ 0  0  0 -1]
      \begin{equation*}
        \begin{mat}
          1 &2 &0  &2 \\
          2 &4 &1  &0 \\
          0 &0 &-1 &3 \\
          3 &-1 &1 &4
        \end{mat}
        \grstep[-3\rho_1+\rho_4]{-2\rho_1+\rho_2}
        \grstep{\rho_2\leftrightarrow\rho_4}
        \grstep{\rho_3+\rho_4}      
        \begin{mat}
          1 &2 &0  &2 \\
          0 &-7 &1  &-2 \\
          0 &0 &-1 &3 \\
          0 &0 &0 &-1
        \end{mat}
      \end{equation*}
      阶梯形矩阵的对角元之积为 
       $1\cdot(-7)\cdot(-1)\cdot(-1)=-7$。
      在高斯消元的过程中没有行被倍乘，但有
      一次行对换，所以求行列式时须变号，
      得 $+7$。
    \end{answer}
  \item 
     证明此式。
      \begin{equation*}
         \det(
         \begin{mat}
             1    &1   &1    \\
             a    &b   &c    \\
             a^2  &b^2 &c^2
         \end{mat}
         )
         =(b-a)(c-a)(c-b)
      \end{equation*}
     \begin{answer}
       利用 $\nbyn{3}$ 矩阵的行列式公式，
       我们将左边展开
       \begin{equation*}
         1\cdot b\cdot c^2+1\cdot c\cdot a^2+1\cdot a\cdot b^2
          -b^2\cdot c\cdot 1 -c^2\cdot a\cdot 1-a^2\cdot b\cdot 1
       \end{equation*}
       而将右边展开则得
       \begin{equation*}
         (bc-ba-ac+a^2)\cdot(c-b)
         =c^2b-b^2c-bac+b^2a-ac^2+acb+a^2c-a^2b
       \end{equation*}
       于是只需验证两者相等即可。
       （\textit{注}。
       这是
       \definend{Vandermonde 行列式}\index{Vandermonde!determinant}%
       \index{determinant!Vandermonde}的 \( \nbyn{3} \) 情形，
       它出现在许多应用中）。
     \end{answer}
   \recommended \item 
      哪些实数 \( x \) 使这个矩阵是奇异的？
      \begin{equation*}
         \begin{mat}
            12-x  &4  \\
            8    &8-x
         \end{mat}
      \end{equation*}
      \begin{answer}
         这个方程
         \begin{equation*}
           0=
           \det(
             \begin{mat}
                12-x  &4  \\
                8    &8-x
             \end{mat}
           )
           =64-20x+x^2
           =(x-16)(x-4) 
       \end{equation*}
       有根 \( x=16 \) 和 \( x=4 \)。  
     \end{answer}
  \item \label{exer:ThreeByThreeDetForm} 
    做高斯消元，以验证本节引言中
    所述的 $\nbyn{3}$ 矩阵公式。
    \begin{center}
      \( \begin{mat}
               a  &b  &c  \\
               d  &e  &f  \\
               g  &h  &i
         \end{mat} \)
      是非奇异的当且仅当
      \( aei+bfg+cdh-hfa-idb-gec \neq 0 \)
    \end{center}
    \begin{answer}
      我们先把矩阵化成阶梯形。
      首先，假设 \( a\neq 0 \) 且 \( ae-bd\neq 0 \)。
      \begin{multline*}
        \grstep{(1/a)\rho_1}
        \begin{mat}
           1   &b/a   &c/a   \\
           d   &e     &f     \\
           g   &h     &i
         \end{mat}                                           
        \grstep[-g\rho_1+\rho_3]{-d\rho_1+\rho_2}
        \begin{mat}
           1   &b/a           &c/a           \\
           0   &(ae-bd)/a     &(af-cd)/a     \\
           0   &(ah-bg)/a     &(ai-cg)/a
         \end{mat}                                            \\
        \grstep{(a/(ae-bd))\rho_2}
        \begin{mat}
           1   &b/a           &c/a             \\
           0   &1             &(af-cd)/(ae-bd) \\
           0   &(ah-bg)/a     &(ai-cg)/a
         \end{mat}
      \end{multline*}
      这一步完成了计算。
      \begin{equation*}
        \grstep{((ah-bg)/a)\rho_2+\rho_3}
        \begin{mat}
           1   &b/a    &c/a             \\
           0   &1      &(af-cd)/(ae-bd)      \\
           0   &0      &(aei+bgf+cdh-hfa-idb-gec)/(ae-bd)
         \end{mat}
      \end{equation*}
      现在假设 $a\neq 0$ 且 \( ae-bd\neq 0 \)，
      原矩阵是非奇异的
      当且仅当上面第 \( 3,3 \) 处的元素非零。
      也就是说，在这些假设下，原矩阵
      是非奇异的当且仅当 $aei+bgf+cdh-hfa-idb-gec\neq 0$，
      即为所证。

      最后我们逐一看一下，若为方便起见而在前一段中所作的假设不成立，
      会发生什么。
      首先，若 \( a\neq 0 \) 但 \( ae-bd=0 \)，那么我们可以对换
      \begin{equation*}
        \begin{mat}
           1   &b/a           &c/a           \\
           0   &0             &(af-cd)/a     \\
           0   &(ah-bg)/a     &(ai-cg)/a
         \end{mat}                                  
        \grstep{\rho_2\leftrightarrow\rho_3}
        \begin{mat}
           1   &b/a           &c/a           \\
           0   &(ah-bg)/a     &(ai-cg)/a     \\
           0   &0             &(af-cd)/a
         \end{mat}
      \end{equation*}
      从而得出结论：矩阵是非奇异的当且仅当
      \( ah-bg=0 \) 或 \( af-cd=0 \)。
      条件‘\( ah-bg=0 \) 或 \( af-cd=0 \)’等价于
      条件‘\( (ah-bg)(af-cd)=0 \)’。
      乘开并利用本情形的假设 $ae-bd=0$，以 $ae$ 代 $bd$，便得
      \begin{multline*}
         0=ahaf-ahcd-bgaf+bgcd
          =ahaf-ahcd-bgaf+aegc            \\
          =a(haf-hcd-bgf+egc)
      \end{multline*}
      由于 \( a\neq 0 \)，故矩阵
      是非奇异的当且仅当 \( haf-hcd-bgf+egc=0 \)。
      因此，在 \( a\neq 0 \) 且 \( ae-bd=0 \) 的情形下，
      当
      \( haf-hcd-bgf+egc-i(ae-bd)=0 \) 时矩阵
      是非奇异的。

      余下的情形是常规的。
      对于 \( a=0 \) 但 \( d\neq 0 \) 的情形，以及 \( a=0 \)、\( d=0 \)
      但 \( g\neq 0 \) 的情形，都先对换两行，然后
      照上面那样继续即可。
      \( a=0 \)、\( d=0 \) 且 \( g=0 \) 的情形则很简单\Dash 该矩阵是
      奇异的，因为它的列构成一个线性相关的组，而且
      行列式算出来为零。  
    \end{answer}
  \item 
     证明 \( \Re^2 \) 中过 \( (x_1,y_1) \)
     与 \( (x_2,y_2) \) 的直线方程由这个行列式给出。
     \begin{equation*}
        \det(
        \begin{mat}
           x   &y   &1  \\
           x_1 &y_1 &1  \\
           x_2 &y_2 &1
        \end{mat})=0 \qquad x_1\neq x_2
     \end{equation*}
     \begin{answer}
       算出行列式并做一点代数运算，便得到
       \begin{align*}
          0
          &=y_1x+x_2y+x_1y_2-y_2x-x_1y-x_2y_1     \\
          (x_2-x_1)\cdot y
          &=(y_2-y_1)\cdot x+x_2y_1-x_1y_2              \\
          y
          &=\frac{y_2-y_1}{x_2-x_1}\cdot x+\frac{x_2y_1-x_1y_2}{x_2-x_1}
       \end{align*}
       注意这是一个直线的方程（特别地，
       其中含有熟悉的斜率表达式），
       并注意 \( (x_1,y_1) \) 与 \( (x_2,y_2) \) 都满足它。 
     \end{answer}
  \item
    许多人学过
    计算这个 \( \nbyn{3} \)
    矩阵行列式的口诀：~把前两列抄到矩阵右侧，
    然后取沿正向对角线的乘积并相加，
    再取沿反向对角线的乘积并相减。
    也就是说，先写成
    \begin{equation*}
        \begin{pmat}{ccc|cc}
          h_{1,1} &h_{1,2} &h_{1,3} &h_{1,1} &h_{1,2} \\
          h_{2,1} &h_{2,2} &h_{2,3} &h_{2,1} &h_{2,2} \\
          h_{3,1} &h_{3,2} &h_{3,3} &h_{3,1} &h_{3,2}
        \end{pmat}
     \end{equation*}
     然后计算这个式子。
     \begin{equation*}
      \begin{array}{l}
      h_{1,1}h_{2,2}h_{3,3}+h_{1,2}h_{2,3}h_{3,1}+h_{1,3}h_{2,1}h_{3,2} \\
      \>-h_{3,1}h_{2,2}h_{1,3}-h_{3,2}h_{2,3}h_{1,1}
        -h_{3,3}h_{2,1}h_{1,2}
      \end{array}
    \end{equation*}
    \begin{exparts}
      \partsitem 验证这与本节引言中
        给出的公式一致。
      \partsitem 它能否推广到其他阶数的行列式？
    \end{exparts}
    \begin{answer}
      \begin{exparts}
        \partsitem 与本节引言中给出的公式比较，很容易。
        \partsitem 它对 \( \nbyn{2} \) 矩阵成立
          \begin{align*}
            \begin{pmat}{cc|c}
              h_{1,1} &h_{1,2} &h_{1,1} \\
              h_{2,1} &h_{2,2} &h_{2,1}
            \end{pmat}
            &=\begin{array}[t]{@{}l@{}}
              h_{1,1}h_{2,2}+h_{1,2}h_{2,1}  \\
              \>-h_{2,1}h_{1,2}-h_{2,2}h_{1,1}
            \end{array}                                     \\
            &=h_{1,1}h_{2,2}-h_{1,2}h_{2,1}
          \end{align*}
          但对 \( \nbyn{4} \) 矩阵不成立。
          例如，这个矩阵是
          奇异的，因为第二行与第三行相等
          \begin{equation*}
            \begin{mat}[r]  
              1  &0  &0  &1   \\
              0  &1  &1  &0   \\
              0  &1  &1  &0   \\
             -1  &0  &0  &1  
            \end{mat}
          \end{equation*}
          但按口诀的方案算并不得到零。
          \begin{equation*}
            \begin{pmat}{rrrr|rrr}
                       1  &0  &0  &1  &1  &0  &0  \\
                       0  &1  &1  &0  &0  &1  &1  \\
                       0  &1  &1  &0  &0  &1  &1  \\
                      -1  &0  &0  &1  &-1 &0  &0
            \end{pmat}
            =\begin{array}[t]{@{}l@{}}
               1+0+0+0 \\
               \>-(-1)-0-0-0
              \end{array}
           \end{equation*}
      \end{exparts}  
     \end{answer}
  \item 
    向量
    \begin{equation*}
      \vec{x}=\colvec{x_1 \\ x_2 \\ x_3}
      \qquad
      \vec{y}=\colvec{y_1 \\ y_2 \\ y_3}
    \end{equation*}
    的
    \definend{叉积}\index{cross product}\index{vector!cross product} 
    是按这个行列式算出的向量。
    \begin{equation*}
      \vec{x}\times\vec{y}=
      \det(\begin{mat}
        \vec{e}_1  &\vec{e}_2  &\vec{e}_3  \\
        x_1        &x_2        &x_3        \\
        y_1        &y_2        &y_3
      \end{mat})
    \end{equation*}
    注意第一行的元素是向量，即 $\Re^3$ 的
    标准基中的向量。
    证明两个向量的叉积与每个向量都垂直。
    \begin{answer}
      该行列式为 
      $
        (x_2y_3-x_3y_2)\vec{e}_1
          +(x_3y_1-x_1y_3)\vec{e}_2
          +(x_1y_2-x_2y_1)\vec{e}_3
      $。
      要检验垂直性，我们检验它与第一个向量的点积
      为零
      \begin{equation*}
        \colvec{x_1 \\ x_2 \\ x_3}
        \dotprod
        \colvec{x_2y_3-x_3y_2 \\ x_3y_1-x_1y_3 \\ x_1y_2-x_2y_1}
        =x_1x_2y_3-x_1x_3y_2+x_2x_3y_1-x_1x_2y_3+x_1x_3y_2-x_2x_3y_1=0
      \end{equation*}
      而它与第二个向量的点积也为零。
      \begin{equation*}
        \colvec{y_1 \\ y_2 \\ y_3}
        \dotprod
        \colvec{x_2y_3-x_3y_2 \\ x_3y_1-x_1y_3 \\ x_1y_2-x_2y_1}
        =x_2y_1y_3-x_3y_1y_2+x_3y_1y_2-x_1y_2y_3+x_1y_2y_3-x_2y_1y_3=0
      \end{equation*}  
    \end{answer}
  \item 
    证明下列每个论断对 $\nbyn{2}$ 矩阵成立。  
    \begin{exparts}
      \partsitem 乘积的行列式
        等于行列式的乘积
        $\det(ST)=\det(S)\cdot\det(T)$。
      \partsitem 若 \( T \) 是可逆的，则
        逆矩阵的行列式等于行列式的倒数
        \( \det(T^{-1})=(\,\det(T)\,)^{-1} \)。
    \end{exparts}
    矩阵 $T$ 与 $T^\prime$ 称为
    \definend{相似的}\index{similar}，若存在
    非奇异矩阵 $P$ 使得 $T^\prime=PTP^{-1}$。
    （我们将在第五章考察这一关系。）
    证明相似的 \( \nbyn{2} \) 矩阵有相同的
    行列式。
    \begin{answer}
       \begin{exparts}
        \partsitem 直接代入计算：
          乘积的行列式为
          \begin{align*}
             \det(\begin{mat}
                       a  &b  \\
                       c  &d
                    \end{mat}
                    \begin{mat}
                       w  &x  \\
                       y  &z
                    \end{mat}  )
             &=
             \det(\begin{mat}
                 aw+by  &ax+bz  \\
                 cw+dy  &cx+dz
              \end{mat} )                 \\
             &=
             \begin{array}[t]{@{}l@{}} 
                 acwx+adwz+bcxy+bdyz  \\
                 \> -acwx-bcwz-adxy-bdyz
             \end{array}
          \end{align*}
          而行列式的乘积为
          \begin{equation*}
             \det(\begin{mat}
                a  &b  \\
                c  &d
             \end{mat})
             \cdot\det(\begin{mat}
                w  &x  \\
                y  &z
             \end{mat})
             =
             (ad-bc)\cdot (wz-xy)
          \end{equation*}
          验证两者相等很容易。
        \partsitem 利用前一小题。
       \end{exparts}  
       \noindent 由上面两条立即得到：相似矩阵有相同的行列式
       $
          \det(PTP^{-1})=\det(P)\cdot\det(T)\cdot\det(P^{-1})
       $。
     \end{answer}
  \recommended \item
    证明平面内这个区域的面积
    \begin{center}
      \includegraphics{det/mp/ch4.30}
    \end{center}
    等于这个行列式的值。
    \begin{equation*}
       \det(
       \begin{mat}
          x_1  &x_2  \\
          y_1  &y_2
       \end{mat})
    \end{equation*}
    并与这个比较。
    \begin{equation*}
       \det(
       \begin{mat}
          x_2  &x_1  \\
          y_2  &y_1
       \end{mat})
    \end{equation*}
    \begin{answer}
      一种方法是数出这些面积
      \begin{center}
        \includegraphics{det/mp/ch4.31}
      \end{center}
      即取整个矩形的面积，再减去
      左上矩形 $A$、上中三角形 $B$、
      右上三角形 $D$、左下三角形 $C$、
      下中三角形 $E$ 与右下矩形 $F$ 的面积
      \( (x_1+x_2)(y_1+y_2)-x_2y_1-(1/2)x_1y_1-(1/2)x_2y_2
              -(1/2)x_2y_2-(1/2)x_1y_1-x_2y_1 \)。
      化简即得行列式公式。

      这个行列式是上面那个的相反数；该公式
      区分了第二列是否在第一列的
      逆时针方向。  
     \end{answer}
  \item 
    证明对 \( \nbyn{2} \) 矩阵，矩阵的行列式
    等于其转置的行列式。
    这对 \( \nbyn{3} \) 矩阵也成立吗？
    \begin{answer}
      对 \( \nbyn{2} \) 矩阵，利用
      引言中引用的公式，计算很容易。
      它对 \( \nbyn{3} \) 矩阵也成立；其
      计算是常规的。  
    \end{answer}
  \item 
    行列式函数是线性的吗 \Dash  即
    \( \det(x\cdot T+y\cdot S)=x\cdot \det(T)+y\cdot \det(S) \) 成立吗？
    \begin{answer}
      不成立。
      我们用 $\nbyn{2}$ 行列式来说明。
      回想常数可以每次提出一行。
      \begin{equation*}
         \det(
         \begin{mat}[r]
            2  &4  \\
            2  &6  \\
         \end{mat})
         =
         2\cdot\det(\begin{mat}[r]
            1  &2  \\
            2  &6  \\
         \end{mat})
         =
         2\cdot 2\cdot \det(\begin{mat}[r]
            1  &2  \\
            1  &3  \\
         \end{mat})
      \end{equation*}
      这与线性性矛盾（这里我们不需要 \( S \)，也就是说，我们可以
      取 $S$ 为零矩阵）。  
    \end{answer}
  \item 
    证明：若 \( A \) 是 \( \nbyn{3} \) 的，则
    \( \det(c\cdot A)=c^3\cdot \det(A) \) 对任意标量 \( c \) 成立。
    \begin{answer}
       把 \( c \) 一次一行地提出来。  
    \end{answer}
  \item 
    哪些实数 \( \theta \) 使
    \begin{equation*}
       \begin{mat}
          \cos\theta  &-\sin\theta  \\
          \sin\theta  &\cos\theta
       \end{mat}
    \end{equation*}
    奇异？
    从几何上解释。
    \begin{answer}
      不存在使该矩阵奇异的实数 \( \theta \)，
      因为矩阵的行列式
      \( \cos^2\theta+\sin^2\theta \) 永不为 $0$，它对所有 $\theta$
      都等于 $1$。
      从几何上看，关于标准基，
      这个矩阵表示
      平面旋转角度 \( \theta \) 的转动。
      每个这样的映射都是单射 \Dash  例如，它是可逆的。  
    \end{answer}
  \puzzle \item  
    \cite{Monthly55p257}
    若一个三阶行列式的元素为
    \( 1 \), \( 2 \), \ldots, \( 9 \)，它可能取到的最大值
    是多少？
    \begin{answer}
      \answerasgiven
      设 \( P \) 为行列式中三个正项之和，
      \( -N \) 为三个负项之和。
      \( P \) 的最大值为
      \begin{equation*}
        9\cdot 8\cdot 7 +6\cdot 5\cdot 4 +3\cdot 2\cdot 1=630.
      \end{equation*}
      与 \( P \) 相容的 \( N \) 的最小值为
      \begin{equation*}
        9\cdot 6\cdot 1 +8\cdot 5\cdot 2 +7\cdot 4\cdot 3=218.
      \end{equation*}
      \( P \) 的任何改变都会使该和的降低量超过
      \( 4 \)。
      因此 \( 412 \) 是该行列式的最大值，而该行列式的一种形式是
      \begin{equation*}
         \begin{vmat}[r]
            9  &4  &2  \\
            3  &8  &6  \\
            5  &1  &7
         \end{vmat}.
      \end{equation*}  
    \end{answer}
\end{exercises}




















\subsection{行列式的性质}
\index{determinant|(}
我们想要一个公式，
用以判断一个 $\nbyn{n}$ 矩阵是否非奇异的。
我们不打算一开始就陈述这样一个公式，
而是先针对每个~$n$，
考虑这样的公式所计算的函数。
我们将通过一列性质来定义这个函数。
随后我们将证明具有这些性质的函数存在且唯一，
并说明如何计算它。
（由于我们最终会证明这一点，所以从一开始
我们就直接说“\( \det(T) \)”而不说
“若存在唯一的行列式函数，则 \( \det(T) \)”。）

\begin{definition}  \label{def:Det}
%<*df:Det>
一个 \( \nbyn{n} \) \definend{行列式\/}\index{determinant!definition}%
\index{matrix!determinant}是一个函数
\( \map{\det}{\matspace_{\nbyn{n}}}{\Re} \)，满足下列条件：
\begin{enumerate}
  \item
    $
       \det (\vec{\rho}_1,\dots,k\cdot\vec{\rho}_i
                                     + \vec{\rho}_j,\dots,\vec{\rho}_n)
       =\det (\vec{\rho}_1,\dots,\vec{\rho}_j,\dots,\vec{\rho}_n)
    $ 当 \( i\ne j \) 时
  \item
    $
       \det (\vec{\rho}_1,\ldots,\vec{\rho}_j,
               \dots,\vec{\rho}_i,\dots,\vec{\rho}_n)
       = -\det (\vec{\rho}_1,\dots,\vec{\rho}_i,\dots,\vec{\rho}_j,
                \dots,\vec{\rho}_n)
    $
    当 \( i\ne j\) 时
  \item
    $
       \det (\vec{\rho}_1,\dots,k\vec{\rho}_i,\dots,\vec{\rho}_n)
       = k\cdot \det (\vec{\rho}_1,\dots,\vec{\rho}_i,\dots,\vec{\rho}_n)
    $
    对任意标量~$k$
    % for \( k\ne 0\)
  \item
     $
       \det(I)=1
     $
     其中 \( I \) 是单位矩阵
\end{enumerate}
（各个 $\vec{\rho}\,$ 是矩阵的行）。
我们常把 \( \det (T) \) 写成 \( \deter{T} \)。
%</df:Det>
\end{definition}

\begin{remark}  \label{rem:SwapRowsRedun}
%<*re:SwapRowsRedun>
条件~(2) 是多余的，因为
 \begin{equation*}
   T\grstep{\rho_i+\rho_j}
    \repeatedgrstep{-\rho_j+\rho_i}
    \repeatedgrstep{\rho_i+\rho_j}
    \repeatedgrstep{-\rho_i}
    \hat{T}
\end{equation*}
对换第 \( i \) 行与第~\( j \) 行。
我们列出它是为了
与前面各章中高斯消元法的表述保持一致。
%</re:SwapRowsRedun>
\end{remark}

\begin{remark}
条件~(3) 没有 $k\neq 0$ 的限制，虽然
高斯消元法中把一行乘以~$k$ 的操作
确有此限制。
% \nearbylemma{le:IdenRowsDetZero} below
下面的结果表明，在这里我们不需要该限制。
\end{remark}

% % \begin{remark}
% The previous subsection's plan
% asks for the determinant of an echelon form matrix to be the product
% down the diagonal.
% The next result shows that in the presence of the other
% three, condition~(4) gives that.
% % \end{remark}

\begin{lemma}   \label{le:IdenRowsDetZero}
%<*lm:IdenRowsDetZero>
有两行相同的矩阵，其行列式为零。
有一个零行的矩阵，其行列式为零。
矩阵是非奇异的当且仅当其行列式非零。
阶梯形矩阵的行列式等于沿其对角线的乘积。
%</lm:IdenRowsDetZero>
\end{lemma}

\begin{proof}
%<*pf:IdenRowsDetZero0>
为验证第一句，对换这两个相同的行。
行列式的符号改变了，但矩阵是同一个，
因而其行列式也相同。
于是行列式为零。
%</pf:IdenRowsDetZero0>

%<*pf:IdenRowsDetZero1>
对于第二句，
将零行乘以二。
这使行列式加倍，但该行
保持不变，因而
行列式也保持不变。
所以行列式必为零。
%</pf:IdenRowsDetZero1>

%<*pf:IdenRowsDetZero2>
对于第三句，
作高斯–若尔当消元
$T \rightarrow\cdots\rightarrow\hat{T}$。
由前三个性质，
$T$ 的行列式为零当且仅当
$\hat{T}$ 的行列式为零
（尽管二者可能在符号或大小上不同）。
非奇异矩阵 $T$ 经高斯–若尔当消元化为单位矩阵，
因而其行列式非零。
奇异的 $T$ 则化为带有一个零行的 $\hat{T}$；
由本引理的第二句，其行列式为零。
%</pf:IdenRowsDetZero2>

%<*pf:IdenRowsDetZero3>
第四句有两种情形。
若阶梯形矩阵是奇异的，则
它有一个零行。
于是其对角线上有一个零，并且
沿其对角线的乘积为零。
由本结论的第三句，行列式为零，
因此该矩阵的行列式等于
沿其对角线的乘积。
%</pf:IdenRowsDetZero3>

%<*pf:IdenRowsDetZero4>
若阶梯形矩阵是非奇异的，则它的对角元
都不为零。
这意味着我们可以除以这些对角元，
再利用条件~(3) 使对角线上都得到 $1$。
\begin{equation*}
  \begin{vmat}
    t_{1,1}  &t_{1,2}  &     &t_{1,n}  \\
    0        &t_{2,2}  &     &t_{2,n}  \\
             &         &\ddots         \\
    0        &         &     &t_{n,n}
  \end{vmat}
  =
  t_{1,1}\cdot t_{2,2}\cdots t_{n,n}\cdot
  \begin{vmat}
    1        &t_{1,2}/t_{1,1}  &     &t_{1,n}/t_{1,1}  \\
    0        &1                &     &t_{2,n}/t_{2,2}  \\
             &                 &\ddots         \\
    0        &                 &     &1
  \end{vmat}
\end{equation*}
然后，高斯–若尔当消元的若尔当部分使之成为单位矩阵。
\begin{equation*}
  =
  t_{1,1}\cdot t_{2,2}\cdots t_{n,n}\cdot
  \begin{vmat}
    1        &0                &     &0                \\
    0        &1                &     &0                \\
             &                 &\ddots         \\
    0        &                 &     &1
  \end{vmat}
  =
  t_{1,1}\cdot t_{2,2}\cdots t_{n,n}\cdot 1
\end{equation*}
所以在这种情形，行列式
同样是沿对角线的乘积。
%</pf:IdenRowsDetZero4>
\end{proof}

这给了我们一种办法，用来计算行列式函数
在一个矩阵上的取值：
作高斯消元，记录行对换所引起的任何变号
以及我们提出的任何标量倍数，
最后沿阶梯形结果的
对角线相乘。
这一算法与高斯消元法一样快，因此
对我们将会看到的
所有矩阵都是实用的。

\begin{example}
用高斯消元法
计算 \( \nbyn{2} \) 行列式
\begin{equation*}
   \begin{vmat}[r]
      2  &4  \\
      -1 &3
   \end{vmat}
   =
   \begin{vmat}[r]
      2  &4  \\
      0  &5
   \end{vmat}
   =10
\end{equation*}
省不了多少时间，
因为 $\nbyn{2}$ 行列式公式很容易。
然而，一个 \( \nbyn{3} \) 行列式
用高斯消元法来计算往往比用它的公式更容易。
\begin{equation*}
   \begin{vmat}[r]
     2  &2  &6  \\
     4  &4  &3  \\
     0  &-3 &5
   \end{vmat}
   =
   \begin{vmat}[r]
     2  &2  &6  \\
     0  &0  &-9 \\
     0  &-3 &5
   \end{vmat}
   =
   -\begin{vmat}[r]
     2  &2  &6  \\
     0  &-3 &5  \\
     0  &0  &-9
   \end{vmat}
   =-54
\end{equation*}
\end{example}

\begin{example}
比 $\nbyn{3}$ 更大的行列式
用高斯消元法的过程很快就完成了。
\begin{equation*}
   \begin{vmat}[r]
      1  &0  &1  &3  \\
      0  &1  &1  &4  \\
      0  &0  &0  &5  \\
      0  &1  &0  &1
   \end{vmat}
   =
   \begin{vmat}[r]
      1  &0  &1  &3  \\
      0  &1  &1  &4  \\
      0  &0  &0  &5  \\
      0  &0  &-1 &-3
   \end{vmat}
   =
   -\begin{vmat}[r]
      1  &0  &1  &3  \\
      0  &1  &1  &4  \\
      0  &0  &-1 &-3 \\
      0  &0  &0  &5
   \end{vmat}
   =-(-5)=5
\end{equation*}
\end{example}

上面的例题引出了一个重要的观点。
本章的引言给出了 $\nbyn{2}$
和 $\nbyn{3}$ 行列式的公式，因此我们知道它们存在，
但对于 $\nbyn{4}$ 或更大的矩阵上的行列式函数却没有给出公式。
作为替代，\nearbydefinition{def:Det} 给出了一个行列式函数
应当具有的性质，
并由此引出用高斯消元法来计算行列式。

然而，对任意一个矩阵，我们都可以用高斯消元法以
多种方式将其化为阶梯形。
例如，给定一个消元过程，我们可以插入一个第一步
把最上面一行乘以~$2$，然后插入一个第二步
把它乘以~$1/2$，从而改变这个过程。
因此我们必须担心，两种不同的高斯消元法消元
可能导致行列式的两个不同的计算值。

也就是说，我们必须验证 \nearbydefinition{def:Det} 给出的是
一个良定义的函数。
接下来的两小节将完成这件事，证明存在
满足该定义的
一个良定义的函数。

但首先我们证明，若存在这样的函数，
则至多只有一个。
上面的例题说明了这个想法：~我们
按照定义的性质得到了~$5$。
因此，虽然我们尚未证明 $\det_{\nbyn{4}}$ 存在，
即存在具有性质 (1)\,--\,(4) 的函数，
但若满足这些性质的函数确实存在，
则我们知道它在上述矩阵上给出的值。

\begin{lemma} \label{lm:DetFcnIsUnique}
%<*lm:DetFcnIsUnique>
对每个 $n$，
若存在 $\nbyn{n}$ 行列式函数，则它是唯一的。
%</lm:DetFcnIsUnique>
\end{lemma}

\begin{proof}
%<*pf:DetFcnIsUnique>
假设有两个函数
$\map{\det_1,\det_2}{\matspace_{\nbyn{n}}}{\Re}$
满足
\nearbydefinition{def:Det} 的性质以及
其推论 \nearbylemma{le:IdenRowsDetZero}。
给定一个方阵~$M$，固定一种实施高斯消元法
将该矩阵化为阶梯形的方式（有多种方式并不要紧，
只需固定其中一种）。
通过像上面的例题中那样
使用这一固定的消元过程\Dash
记录行倍乘的因子以及符号在行对换时如何变号，
然后沿阶梯形结果的对角线相乘\Dash
我们可以算出这两个函数在~$M$ 上必定返回的值，
并且它们必定返回相同的值。
由于它们在每个输入上都给出相同的输出，所以它们是同一个函数。
%</pf:DetFcnIsUnique>
\end{proof}

“若存在 $\nbyn{n}$ 行列式函数”
这一说法强调，
虽然我们可以
用高斯消元法计算出行列式函数
所可能返回的唯一值，
但我们尚未证明这样的函数对一切 $n$ 都存在。
本节其余部分将完成这件事。



\begin{exercises}
  \item[{\em 做这些题时，假设 $\nbyn{n}$ 行列式
      函数对一切 $n$ 都存在。}]
  \recommended \item 通过执行一次行变换求每个行列式。
    \begin{exparts*}
      \partsitem \( \begin{vmat}[r]
                 1  &-2  &1  &2 \\
                 2  &-4  &1  &0 \\
                 0  &0  &-1  &0 \\
                 0  &0  &0  &5
               \end{vmat} \)
      \partsitem \( \begin{vmat}[r]
                 1  &1  &-2  \\
                 0  &0  &4  \\
                 0  &3  &-6
               \end{vmat} \)
    \end{exparts*}
    \begin{answer}
      \begin{exparts}
       \partsitem 作 $2\rho_1+\rho_2$ 化为阶梯形，然后沿对角线相乘。
          行列式为~$0$。
          % sage: M = matrix(QQ, [[1,-2,1,2], [2,-4,1,0], [0,0,-1,0], [0,0,0,5]])
          % sage: M.determinant()
          % 0
        \partsitem 对换第二行与第三行即可将其化为阶梯形（同时改变行列式的符号）。
          沿对角线相乘得~$12$，故所给矩阵的行列式为~$-12$。
          % sage: M = matrix(QQ, [[1,1,-2], [0,0,4], [0,3,-6]])
          % sage: M.determinant()
          % -12
      \end{exparts}
    \end{answer}
  \recommended \item 
    用高斯消元法求每个行列式。
    \begin{exparts*}
      \partsitem \( \begin{vmat}[r]
                 3  &1  &2  \\
                 3  &1  &0  \\
                 0  &1  &4
               \end{vmat} \)
      \partsitem \( \begin{vmat}[r]
                 1  &0  &0  &1 \\
                 2  &1  &1  &0 \\
                -1  &0  &1  &0 \\
                 1  &1  &1  &0
               \end{vmat} \)
    \end{exparts*}
    \begin{answer}
      \begin{exparts}
      \partsitem \( \begin{vmat}[r]
                 3  &1  &2  \\
                 3  &1  &0  \\
                 0  &1  &4
               \end{vmat}
               =
               \begin{vmat}[r]
                 3  &1  &2  \\
                 0  &0  &-2 \\
                 0  &1  &4
               \end{vmat}
               =
               -\begin{vmat}[r]
                 3  &1  &2  \\
                 0  &1  &4  \\
                 0  &0  &-2
               \end{vmat}
               =6  \)
      \partsitem \( \begin{vmat}[r]
                 1  &0  &0  &1 \\
                 2  &1  &1  &0 \\
                -1  &0  &1  &0 \\
                 1  &1  &1  &0
               \end{vmat}
               =
               \begin{vmat}[r]
                 1  &0  &0  &1 \\
                 0  &1  &1  &-2\\
                 0  &0  &1  &1 \\
                 0  &1  &1  &-1
               \end{vmat}
               =
               \begin{vmat}[r]
                 1  &0  &0  &1 \\
                 0  &1  &1  &-2\\
                 0  &0  &1  &1 \\
                 0  &0  &0  &1
               \end{vmat}
               =1    \)
      \end{exparts} 
    \end{answer}
  \item 用高斯消元法求每个行列式。
    \begin{exparts*}
      \partsitem \( \begin{vmat}[r]
                 2  &-1  \\
                -1 &-1
               \end{vmat}  \)
      \partsitem \( \begin{vmat}[r]
                 1  &1  &0  \\
                 3  &0  &2  \\
                 5  &2  &2
               \end{vmat} \)
    \end{exparts*}
    \begin{answer}
     \begin{exparts}
      \partsitem \(
        \begin{vmat}[r]
                 2  &-1  \\
                 -1 &-1
               \end{vmat}
             =\begin{vmat}[r]
                 2  &-1  \\
                 0  &-3/2
               \end{vmat}=-3  \);
      \partsitem \( \begin{vmat}[r]
                 1  &1  &0  \\
                 3  &0  &2  \\
                 5  &2  &2
               \end{vmat}
              =\begin{vmat}[r]
                 1  &1  &0  \\
                 0  &-3 &2  \\
                 0  &-3 &2
               \end{vmat}
              =\begin{vmat}[r]
                 1  &1  &0  \\
                 0  &-3 &2  \\
                 0  &0  &0
               \end{vmat}
              =0 \)
     \end{exparts}  
    \end{answer}
  \item 
    对哪些 \( k \) 值，该方程组
    有唯一解？
    \begin{equation*}
       \begin{linsys}{4}
          x  &  &  &+ &z  &-  &w  &=  &2  \\
             &  &y &- &2z &   &   &=  &3  \\
          x  &  &  &+ &kz &   &   &=  &4  \\
             &  &  &  &z  &-  &w  &=  &2
        \end{linsys}
    \end{equation*}
    \begin{answer}
     什么时候行列式不为零？
      \begin{equation*}
         \begin{vmat}
           1  &0  &1  &-1  \\
           0  &1  &-2 &0   \\
           1  &0  &k  &0   \\
           0  &0  &1  &-1
         \end{vmat}
         =
         \begin{vmat}
           1  &0  &1  &-1  \\
           0  &1  &-2 &0   \\
           0  &0  &k-1&1   \\
           0  &0  &1  &-1
         \end{vmat}           
      \end{equation*}
      显然，
      $k=1$ 给出非奇异性，从而行列式非零。
      若 $k\neq 1$，则
      用倍加 $(-1/k-1)\rho_3+\rho_4$ 可得到阶梯形。
      \begin{equation*}
         =
         \begin{vmat}
           1  &0  &1  &-1  \\
           0  &1  &-2 &0   \\
           0  &0  &k-1&1   \\
           0  &0  &0  &-1-(1/k-1)
         \end{vmat}
      \end{equation*}
      沿对角线相乘得 $(k-1)(-1-(1/k-1))=-(k-1)-1=-k$。
      因此，该矩阵的行列式非零，从而方程组
      有唯一解，当且仅当 \( k\neq 0 \)。  
    \end{answer}
  \recommended \item 
    将下列各式用 \( \deter{H} \) 表示。
    \begin{exparts}
      \partsitem \( \begin{vmat}
            h_{3,1}  &h_{3,2} &h_{3,3} \\
            h_{2,1}  &h_{2,2} &h_{2,3} \\
            h_{1,1}  &h_{1,2} &h_{1,3}
          \end{vmat} \)
      \partsitem \( \begin{vmat}
           -h_{1,1}   &-h_{1,2}  &-h_{1,3} \\
           -2h_{2,1}  &-2h_{2,2} &-2h_{2,3} \\
           -3h_{3,1}  &-3h_{3,2} &-3h_{3,3}
          \end{vmat} \)
      \partsitem \( \begin{vmat}
            h_{1,1}+h_{3,1}  &h_{1,2}+h_{3,2} &h_{1,3}+h_{3,3} \\
            h_{2,1}          &h_{2,2}         &h_{2,3} \\
            5h_{3,1}         &5h_{3,2}        &5h_{3,3}
          \end{vmat} \)
    \end{exparts}
    \begin{answer}
      \begin{exparts}
        \partsitem 行列式定义中的条件~(2)
          经对换 $\rho_1\leftrightarrow\rho_3$ 适用。
          \begin{equation*}
            \begin{vmat}
              h_{3,1}  &h_{3,2} &h_{3,3} \\
              h_{2,1}  &h_{2,2} &h_{2,3} \\
              h_{1,1}  &h_{1,2} &h_{1,3}
            \end{vmat}
            =
            -\begin{vmat}
              h_{1,1}  &h_{1,2} &h_{1,3} \\
              h_{2,1}  &h_{2,2} &h_{2,3} \\
              h_{3,1}  &h_{3,2} &h_{3,3} 
            \end{vmat}
          \end{equation*}
        \partsitem 条件~(3) 适用。
          \begin{multline*}
            \begin{vmat}
             -h_{1,1}   &-h_{1,2}  &-h_{1,3} \\
             -2h_{2,1}  &-2h_{2,2} &-2h_{2,3} \\
             -3h_{3,1}  &-3h_{3,2} &-3h_{3,3}
            \end{vmat}
            =
            (-1)\cdot(-2)\cdot(-3)\cdot
            \begin{vmat}
             h_{1,1}   &h_{1,2}  &h_{1,3} \\
             h_{2,1}   &h_{2,2}  &h_{2,3} \\
             h_{3,1}   &h_{3,2}  &h_{3,3}
            \end{vmat}                                  \\
            =
            (-6)\cdot
            \begin{vmat}
             h_{1,1}   &h_{1,2}  &h_{1,3} \\
             h_{2,1}   &h_{2,2}  &h_{2,3} \\
             h_{3,1}   &h_{3,2}  &h_{3,3}
            \end{vmat}
          \end{multline*}
        \partsitem 
          \begin{multline*}
          \begin{vmat}
            h_{1,1}+h_{3,1}  &h_{1,2}+h_{3,2} &h_{1,3}+h_{3,3} \\
            h_{2,1}          &h_{2,2}         &h_{2,3} \\
            5h_{3,1}         &5h_{3,2}        &5h_{3,3}
          \end{vmat}                                                 \\
          \begin{aligned}
          &=
          5\cdot     
          \begin{vmat}
            h_{1,1}+h_{3,1}  &h_{1,2}+h_{3,2} &h_{1,3}+h_{3,3} \\
            h_{2,1}          &h_{2,2}         &h_{2,3} \\
            h_{3,1}          &h_{3,2}         &h_{3,3}
          \end{vmat}                                     \\
          &=5\cdot
          \begin{vmat}
            h_{1,1}          &h_{1,2}         &h_{1,3} \\
            h_{2,1}          &h_{2,2}         &h_{2,3} \\
            h_{3,1}          &h_{3,2}         &h_{3,3}
          \end{vmat}
          \end{aligned}
          \end{multline*}
      \end{exparts}  
     \end{answer}
  \recommended \item 
    求对角矩阵的行列式。
    \begin{answer}
       对角矩阵已是阶梯形，所以
       行列式就是对角线上元素的乘积。
    \end{answer}
  \item 
    若系数矩阵的行列式非零，试描述一个齐次线性方程组的解集。
    \begin{answer}
      它就是平凡子空间。  
    \end{answer}
  \recommended \item 
    证明此行列式为零。
    \begin{equation*}
      \begin{vmat}
        y+z  &x+z  &x+y  \\
        x    &y    &z    \\
        1    &1    &1
      \end{vmat}
    \end{equation*}
    \begin{answer} 
      把第二行加到第一行，得到的矩阵其第一行
      是第三行的 $x+y+z$ 倍。  
    \end{answer}
  \item 
    \begin{exparts}
      \partsitem 求以 $(-1)^{i+j}$ 为 $i,j$ 元的
        $\nbyn{1}$、$\nbyn{2}$ 与 $\nbyn{3}$ 矩阵。  
      \partsitem 求以 \( (-1)^{i+j} \) 为 \( i,j \)
        元的方阵的行列式。
    \end{exparts}
    \begin{answer}
      \begin{exparts}
        \partsitem 
          $\begin{mat}[r]
            1
          \end{mat}$,
          $\begin{mat}[r]
            1  &-1  \\
            -1  &1
          \end{mat}$,
          $\begin{mat}[r]
            1   &-1  &1   \\
            -1  &1   &-1  \\
            1   &-1  &1
          \end{mat}$
        \partsitem $\nbyn{1}$ 情形的行列式为 $1$。
          在其余各种情形中，第二行都是第一行的相反数，
          因此矩阵是奇异的，行列式为零。 
      \end{exparts} 
    \end{answer}
  \item 
    \begin{exparts}
      \partsitem 求以 $i+j$ 为 $i,j$ 元的
        $\nbyn{1}$、$\nbyn{2}$ 与 $\nbyn{3}$ 矩阵。  
      \partsitem 求以 $i+j$ 为 
        $i,j$ 元的方阵的行列式。
    \end{exparts}
    \begin{answer}
      \begin{exparts}
        \partsitem 
          $\begin{mat}[r]
            2
          \end{mat}$,
          $\begin{mat}[r]
            2  &3  \\
            3  &4
          \end{mat}$,
          $\begin{mat}[r]
            2  &3  &4  \\
            3  &4  &5  \\
            4  &5  &6
          \end{mat}$
        \partsitem $\nbyn{1}$ 与 $\nbyn{2}$ 情形给出如下结果。
          \begin{equation*}
            \begin{vmat}[r]
              2
            \end{vmat}
            =2
            \qquad
            \begin{vmat}[r]
              2  &3  \\
              3  &4
            \end{vmat}=-1
          \end{equation*}
          而 $n\geq 3$ 的 $\nbyn{n}$ 矩阵都是奇异的，例如
          \begin{equation*}
            \begin{vmat}[r]
              2  &3  &4  \\
              3  &4  &5  \\
              4  &5  &6
            \end{vmat}=0
          \end{equation*}
          因为第二行的两倍减去第一行
          等于第三行。
          验证这一点是例行的。
      \end{exparts} 
    \end{answer}
  \recommended \item 
    通过给出一个 \( \deter{A+B}\neq\deter{A}+\deter{B} \) 的例子，
    证明行列式函数不是线性的。
    \begin{answer}
      取这一个
      \begin{equation*}
         A=B=
         \begin{mat}[r]
           1  &2  \\
           3  &4
         \end{mat}
      \end{equation*}
      容易验证
      \begin{equation*}
         \deter{A+B}
         =
         \begin{vmat}[r]
           2  &4  \\
           6  &8
         \end{vmat}
         =-8
         \qquad
         \deter{A}+\deter{B}
         =
         -2-2
         =-4
      \end{equation*}
      顺便说一下，这也给出了一个标量乘法不被保持的例子：
      $\deter{2\cdot A}\neq 2\cdot\deter{A}$。  
     \end{answer}
  \item 
    定义中的第二个条件，即行对换改变
    行列式的符号，多少有些令人烦恼。
    它意味着我们必须记录对换的次数，以计算
    符号如何交替变化。
    我们能否摆脱它？
    能否把它换成行对换使行列式
    保持不变这一条件？
    （若可行，那么我们就需要新的 $\nbyn{1}$、
     $\nbyn{2}$ 与 $\nbyn{3}$ 公式，但那只是小事一桩。）
    \begin{answer}
      不行，不能替换。
      \nearbyremark{rem:SwapRowsRedun}
      表明替换之后的四个条件会
      相互矛盾 \Dash  没有函数同时满足这四个条件。
    \end{answer}
  \item 
    证明任何三角矩阵的行列式，无论是上三角的还是下三角的，
    都是其对角线上元素的乘积。
    \begin{answer}
      上三角矩阵已是阶梯形。

      下三角矩阵要么奇异要么非奇异。
      若它奇异，则其对角线上有零，从而其
      行列式（即零）确实是对角线上元素的乘积。 
      若它非奇异，则其对角线上没有零，并且
      我们可以用高斯消元法将其化简为阶梯形，
      而不改变对角线。  
    \end{answer}
  \item
    参见“矩阵乘法的机制”小节中初等矩阵的定义。
    \begin{exparts}
      \partsitem 每种初等矩阵的行列式是什么？
      \partsitem 证明若 \( E \) 是任一初等矩阵，则对任意阶数匹配的
        \( S \) 都有
        \( \deter{ES}=\deter{E}\deter{S} \)。
      \partsitem \textit{（此问题不涉及行列式。）}
        证明若 \( T \) 奇异，则乘积 \( TS \)
        也奇异。
      \partsitem 证明 \( \deter{TS}=\deter{T}\deter{S} \)。
      \partsitem 证明若 \( T \) 非奇异，则
          \( \deter{T^{-1}}=\deter{T}^{-1} \)。
    \end{exparts}
    \begin{answer}
      \begin{exparts}
        \partsitem 行列式定义中的诸性质表明
          \( \deter{M_i(k)}=k \)、
          \( \deter{P_{i,j}}=-1 \)
          以及
          \( \deter{C_{i,j}(k)}=1 \)。
        \partsitem 回想左乘每类矩阵的作用，三种情形都容易验证。
        \partsitem 若 \( TS \) 可逆，即 \( (TS)M=I \)，则
          由矩阵乘法的结合律，
          \( T(SM)=I \) 表明 \( T \) 可逆。
          所以若 \( T \) 不可逆，则 \( TS \) 也不可逆。
        \partsitem 若 \( T \) 奇异，则应用前面的答案：
          \( \deter{TS}=0 \) 且
          \( \deter{T}\cdot\deter{S}=0\cdot\deter{S}=0 \)。
          若 \( T \) 不奇异，则可将它写成初等矩阵的乘积
          $
            \deter{TS}=\deter{E_r\cdots E_1S}=
            \deter{E_r}\cdots\deter{E_1}\cdot\deter{S}=
            \deter{E_r\cdots E_1}\deter{S}=\deter{T}\deter{S}
          $。
        \partsitem \( 1=\deter{I}=\deter{T\cdot T^{-1}}
                     =\deter{T}\deter{T^{-1}} \)
      \end{exparts}  
     \end{answer} 
  \item 
    按下面的方式证明乘积的行列式等于行列式的乘积
    \( \deter{TS}=\deter{T}\,\deter{S} \)。
    取定 \( \nbyn{n} \) 矩阵 \( S \)，考虑函数
    \( \map{d}{\matspace_{\nbyn{n}}}{\Re} \)，其由
    \( T\mapsto \deter{TS}/\deter{S} \) 给出。
    \begin{exparts}
      \partsitem 验证 \( d \) 满足行列式函数定义中的条件~(1)。
      \partsitem 验证条件~(2)。
      \partsitem 验证条件~(3)。
      \partsitem 验证条件~(4)。
      \partsitem 由此得出结论：乘积的行列式等于行列式的乘积。
    \end{exparts}
    \begin{answer}
      \begin{exparts}
        \partsitem 我们必须证明：若
          \begin{equation*}
            T\grstep{k\rho_i+\rho_j}\hat{T}
          \end{equation*}
          则 $d(T)=\deter{TS}/\deter{S}
          =\deter{\hat{T}S}/\deter{S}=d(\hat{T})$。
          如果我们能证明先组合行、再相乘得到 \( \hat{T}S \)，
          与先相乘得到 \( TS \)、再组合行，结果相同
          （因为行列式 \( \deter{TS} \) 不受该倍加的影响，所以我们就有
          \( \deter{\hat{T}S}=\deter{TS} \)，从而 \( d(\hat{T})=d(T) \)），
          便完成证明。
          该论证如下：~把 \( TS \) 的第~\( i \) 行的 \( k \) 倍加到
          \( TS \) 的第~$j$ 行之后，第 \( j,p \) 元为
          \( (kt_{i,1}+t_{j,1})s_{1,p}+\dots+(kt_{i,r}+t_{j,r})s_{r,p} \)，
          这正是 \( \hat{T}S \) 的第 \( j,p \) 元。
        \partsitem 我们只需证明：对换
          $
            T\smash[b]{\grstep{\rho_i\swap\rho_j}}\hat{T}
          $
          之后相乘得到 \( \hat{T}S \)，与先把 \( T \) 乘以 \( S \)、再对换行，
          结果相同（因为行列式 \( \deter{TS} \) 在行对换下变号，
         所以我们就有 \( \deter{\hat{T}S}=-\deter{TS} \)，
          于是 \( d(\hat{T})=-d(T) \)）。
          该论证与前一个论证的进行方式相同。
        \partsitem 正如现在所期待的，我们只需证明：
          把某行乘以一个标量
          $
            T\smash[b]{\grstep{k\rho_i}}\hat{T}
          $
          之后计算 \( \hat{T}S \)，与先计算 \( TS \)、再把该行乘以 \( k \)，
          结果相同（由于行列式 \( \deter{TS} \) 在该倍乘下被乘以 \( k \)，
          我们就有 \( \deter{\hat{T}S}=k\deter{TS} \)，
          所以 \( d(\hat{T})=k\,d(T) \)）。
          论证的进行方式与上面完全一样。
        \partsitem 显然。
        \partsitem 由于我们已经证明 \( d(T) \) 是一个行列式，
          并且行列式函数（若存在）是唯一的，
          我们便有
          所以 \( \deter{T}=d(T)=\deter{TS}/\deter{S} \)。
      \end{exparts}  
    \end{answer}
  \item 
    给定矩阵 $A$ 的\definend{子矩阵}\index{submatrix}\index{matrix!submatrix} 
    指删除 $A$ 的某些行和某些列
    之后得到的矩阵。
    于是，这里第一个矩阵是第二个矩阵的子矩阵。
    \begin{equation*}
      \begin{mat}[r]
        3  &1  \\
        2  &5
      \end{mat}
      \qquad
      \begin{mat}[r]
        3  &4  &1  \\
        0  &9  &-2 \\
        2  &-1 &5
      \end{mat}
    \end{equation*}
    证明：对任意方阵，
    矩阵的秩为 $r$ 当且仅当
    \( r \) 是使得存在一个行列式非零的
    \( \nbyn{r} \) 子矩阵的
    最大整数。
    \begin{answer}
      我们首先论证：秩为 \( r \) 的矩阵有一个行列式非零的
      \( \nbyn{r} \) 子矩阵。
      秩为 \( r \) 的矩阵有一个由 \( r \) 行组成的线性无关集合。
      由这些行构成的矩阵行秩为 \( r \)，因而其
      列秩也为 \( r \)。
      结论：从这 \( r \) 行中可以取出一组由 \( r \) 列组成的线性无关集合，
      于是原矩阵有一个秩为 \( r \) 的
      \( \nbyn{r} \) 子矩阵。

      最后我们证明：若 \( r \) 是具有该性质的最大整数，
      则矩阵的秩为 \( r \)。
      由 \( r \) 的最大性，我们只需证明：
      若一个矩阵有行列式非零的
      \( \nbyn{k} \) 子矩阵，则该矩阵的秩至少为 \( k \)。
      考虑这样一个 \( \nbyn{k} \) 子矩阵。
      它的行是原矩阵的行的一部分，显然
      完整的行的集合是线性无关的。
      于是原矩阵的行秩至少为 \( k \)，而矩阵的行秩
      等于其秩。  
   \end{answer}
  \item
    证明元素全为有理数的矩阵的行列式是有理数。
    \begin{answer}
      元素全为有理数的矩阵用高斯
      消元法化简时，只用到有理数运算即可化为阶梯形矩阵。
      因此对角线上的元素必定是有理数，从而
      对角线上元素的乘积是有理数。  
    \end{answer}
  \puzzle \item 
     \cite{Monthly53p115}
     求以下各组操作的共同点：(a)~约简一个分数，(b)~给鼻子扑粉，
     (c)~给教堂修建新台阶，(d)~让荣休教授留驻校园，
     (e)~把 \( B \)、\( C \)、\( D \) 放入行列式
     \begin{equation*}
       \begin{vmat}
         1   &a   &a^2  &a^3  \\
         a^3 &1   &a    &a^2  \\
         B   &a^3 &1    &a    \\
         C   &D   &a^3  &1
       \end{vmat}.
     \end{equation*}
     \begin{answer}
       \answerasgiven
       行列式的值 \( (1-a^4)^3 \) 与
       \( B \)、\( C \)、\( D \) 的取值无关。
       因此操作~(e) 不改变行列式的值，
       只是改变了它的外观。
       于是在 (a)、(b)、(c)、(d)、(e) 中，
       共同点仅仅在于主要对象的外观被改变了。
       同样的共同点还出现在 (f)~改换玫瑰的名称标签、
       (g)~把一个十进制整数写成 \( 12 \) 进制、(h)~给百合花镀金、
       (i)~给政客粉饰门面，以及 (j)~授予荣誉学位。
    \end{answer}
\end{exercises}




















\subsection{置换展开}
前一小节定义了：满足四个条件的函数就是行列式，并证明了对每个 $n$，
$\nbyn{n}$ 行列式函数至多有一个。
剩下的工作就是证明对每个 $n$ 都存在这样的函数。

但是，我们很容易计算行列式：
用高斯消元法，记录行对换带来的符号变化，
最后沿对角线相乘。
行列式怎么可能不存在呢？

困难在于证明该计算给出的是良定义的\Dash 即唯一的\Dash 结果。
%<*DifferentGaussMethodReductions>
考虑对同一个矩阵进行如下两次高斯消元法化简，
第一次不包含行对换
\begin{equation*}
  \begin{mat}[r]
    1  &2  \\
    3  &4
  \end{mat}
  \grstep{-3\rho_1+\rho_2}
  \begin{mat}[r]
    1  &2  \\
    0  &-2
  \end{mat}
\end{equation*}
第二次包含一次行对换。
\begin{equation*}
  \begin{mat}[r]
    1  &2  \\
    3  &4
  \end{mat}
  \grstep{\rho_1\leftrightarrow\rho_2}
  \begin{mat}[r]
    3  &4  \\
    1  &2
  \end{mat}
  \grstep{-(1/3)\rho_1+\rho_2}
  \begin{mat}[r]
    3  &4  \\
    0  &2/3
  \end{mat}
\end{equation*}
两者都得到行列式 $-2$，
因为在第二种化简中我们记下了：行对换会改变
沿对角线相乘所得结果的符号。
%</DifferentGaussMethodReductions>
我们可以按两种方式进行计算，这一事实引出了两种方式可能给出不同答案的问题。
% To illustrate how a computation that is like the ones that we are doing could 
% fail to be well-defined, 
% suppose that \nearbydefinition{def:Det} did not include condition~(2). 
% That is, suppose that we instead tried to define determinants so that
% the value would not change on a row swap.
% Then first reduction above would
% yield $-2$ while the second would yield $+2$.
% We could still do computations but they wouldn't give consistent outcomes\Dash
% there is no 
% function that satisfies conditions (1), (3), (4), and also this
% altered second condition. 
% Of course, observing that \nearbydefinition{def:Det} does the right thing
% with these two reductions of the above matrix is not enough.
也就是说，我们给出的这种计算行列式值的方法，并没有明显排除如下可能性：比如说，
某个 $\nbyn{7}$~矩阵的两次化简可能导出
不同的行列式值。
若是那样，我们就不再有函数了，因为函数的定义要求每个输入恰好对应唯一的输出。
本节余下的部分将证明，定义~\nearbydefinition{def:Det}
永远不会导致冲突。

% \begin{remark}
% The natural way approach to this would be 
% to define the determinant function with:~``The value of the
% function is the result of doing Gauss's Method,
% keeping track of row swaps, and finishing by multiplying down the diagonal''. 
% (Since Gauss's Method allows for some variation, as we saw above, 
% we would have to fix an explicit algorithm.)
% Then we would be done if we verified that 
% this way algorithm satisfies the four properties.
% However, how to verify this is not evident.
% So the development below will not proceed in this way.
% \end{remark}

为此，我们将定义另一种求行列式值的方法。
（这种替代方法在实践中用处不大，因为它
速度慢。
但它对理论非常有用。）
关键的想法是，
\nearbydefinition{def:Det} 的条件~(3)
表明行列式函数不是线性的。

\begin{example}
利用条件~(3)，标量可以从每一行中分别提出来，
\begin{equation*}
   \begin{vmat}[r]
       4  &2  \\
      -2  &6
     \end{vmat}
  =2\cdot\begin{vmat}[r]
       2  &1  \\
      -2  &6
     \end{vmat}
  =4\cdot\begin{vmat}[r]
       2  &1  \\
      -1  &3
     \end{vmat}
\end{equation*}
而不是从整个矩阵一次提出。
于是，当
\begin{equation*}
   A=\begin{mat}[r]
       2  &1  \\
      -1  &3
     \end{mat}
\end{equation*} 
时，\( \det(2A) \neq 2\cdot\det(A) \) 
（相反，\( \det(2A) =  4\cdot\det(A) \)）。
\end{example}

由于标量是逐行提出的，我们可以猜测
行列式是逐行线性的。

\begin{definition} \label{def:multilinear}
%<*df:Multilinear>
设 \( V \) 是一个向量空间。
映射 \( \map{f}{V^n}{\Re} \) 称为
\definend{多重线性的}\index{function!multilinear}\index{multilinear}，
若  
\begin{enumerate}
  \item 
    $
      f(\vec{\rho}_1,\dots,\vec{v}+\vec{w},
      \ldots,\vec{\rho}_n)
      =f(\vec{\rho}_1,\dots,\vec{v},\dots,\vec{\rho}_n)
      +f(\vec{\rho}_1,\dots,\vec{w},\dots,\vec{\rho}_n)
    $
  \item 
    $
      f(\vec{\rho}_1,\dots,k\vec{v},\dots,\vec{\rho}_n)
      =k\cdot f(\vec{\rho}_1,\dots,\vec{v},\dots,\vec{\rho}_n)
    $
\end{enumerate}
其中 \( \vec{v}, \vec{w}\in V \)，\( k\in\Re \)。
%</df:Multilinear>
\end{definition}

\begin{lemma}  \label{lem:DetsMultilinear}
%<*lm:DetsMultilinear>
行列式是多重线性的。
%</lm:DetsMultilinear>
\end{lemma}

\begin{proof}
%<*pf:DetsMultilinear0>
这里的性质~(2) 恰是 \nearbydefinition{def:Det} 的条件~(3)，
所以我们只需验证性质~(1)。
%</pf:DetsMultilinear0>
% \begin{equation*}
%   \det(\vec{\rho}_1,\dots,\vec{v}+\vec{w},
%       \dots,\vec{\rho}_n)
%   =\det(\vec{\rho}_1,\dots,\vec{v},\dots,\vec{\rho}_n)
%   +\det(\vec{\rho}_1,\dots,\vec{w},\dots,\vec{\rho}_n)
% \end{equation*}

%<*pf:DetsMultilinear1>
分两种情况。
若其余各行构成的集合
\(
  \set{\vec{\rho}_1,\dots,\vec{\rho}_{i-1},\vec{\rho}_{i+1},
       \dots,\vec{\rho}_n} 
\)
是线性相关的，则这三个矩阵都是奇异的，从而三个行列式都为零，等式平凡地成立。
%</pf:DetsMultilinear1>

%<*pf:DetsMultilinear2>
因此假设其余各行构成的集合是线性无关的。
我们可以再添加一个向量
$\sequence{\vec{\rho}_1,\dots,\vec{\rho}_{i-1},\vec{\beta},
               \vec{\rho}_{i+1},\dots,\vec{\rho}_n}$
来作出一个基。
关于这个基表示 $\vec{v}$ 和 $\vec{w}$
\begin{align*}
  \vec{v} &=v_1\vec{\rho}_1+\dots+v_{i-1}\vec{\rho}_{i-1}+v_i\vec{\beta}
            +v_{i+1}\vec{\rho}_{i+1}+\dots+v_n\vec{\rho}_n                \\
  \vec{w} &= w_1\vec{\rho}_1+\dots+w_{i-1}\vec{\rho}_{i-1}+w_i\vec{\beta}
            +w_{i+1}\vec{\rho}_{i+1}+\dots+w_n\vec{\rho}_n
\end{align*}
再相加。
\begin{equation*}
  \vec{v}+\vec{w}
  =
  (v_1+w_1)\vec{\rho}_1+\dots+(v_i+w_i)\vec{\beta}
            +\dots+(v_n+w_n)\vec{\rho}_n 
\end{equation*}
考虑 (1) 的左边并展开 $\vec{v}+\vec{w}$。
% $\det (\vec{\rho}_1,\dots,\vec{v}+\vec{w},\dots,\vec{\rho}_n)$.
\begin{equation*}
    \det (\vec{\rho}_1,\dots,\,(v_1+w_1)\vec{\rho}_1+\dots+(v_i+w_i)\vec{\beta}
            +\dots+(v_n+w_n)\vec{\rho}_n,\,\dots,\vec{\rho}_n) 
   \tag{$*$}
\end{equation*}
% Gauss's Method enables us to eliminate the terms for $\vec{\rho}_1$,
% $\vec{\rho}_2$, etc., from the expanded expression.
% the $(v_1+w_1)\vec{\rho}_1+\dots+(v_i+w_i)\vec{\beta}
%             +\dots+(v_n+w_n)\vec{\rho}_n$ expression.
由行列式定义的条件~(1)，($*$) % the determinant
% $\det(\vec{\rho}_1,\dots,\vec{v}+\vec{w},\dots,\vec{\rho}_n)$
的值在下述操作下不变：
将 $-(v_1+w_1)\vec{\rho}_1$ 倍加到第 $i$ 行 $\vec{v}+\vec{w}$ 上。
第 $i$ 行变成这样。
\begin{equation*}
  \vec{v}+\vec{w}-(v_1+w_1)\vec{\rho}_1
  =
  (v_2+w_2)\vec{\rho}_2+\cdots+(v_i+w_i)\vec{\beta}
            +\dots+(v_n+w_n)\vec{\rho}_n 
\end{equation*}
%</pf:DetsMultilinear2>
%<*pf:DetsMultilinear3>
接着倍加 $-(v_2+w_2)\vec{\rho}_2$ 等等，以消去
来自其他各行的所有项。
应用行列式定义中的条件~(3)。
\begin{multline*}
  \det (\vec{\rho}_1,\dots,\vec{v}+\vec{w},\dots,\vec{\rho}_n)          \\
  \begin{aligned}
    &=\det (\vec{\rho}_1,\dots,(v_i+w_i)\cdot\vec{\beta},\dots,\vec{\rho}_n) \\
    &=(v_i+w_i)\cdot\det (\vec{\rho}_1,\dots,\vec{\beta},\dots,\vec{\rho}_n) \\
    &=v_i\cdot \det (\vec{\rho}_1,\dots,\vec{\beta},\dots,\vec{\rho}_n)  
     +w_i\cdot \det (\vec{\rho}_1,\dots,\vec{\beta},\dots,\vec{\rho}_n)
  \end{aligned}
\end{multline*}
现在这是两个行列式的和。
为完成证明，把 $v_i$ 与 $w_i$ 重新移回到
$\vec{\beta}$ 的前面，
并再次使用行倍加，
这一次用来按基重构 $\vec{v}$
和 $\vec{w}$ 的表达式。
也就是说，从下述操作开始：
将 $v_1\vec{\rho}_1$ 倍加到 $v_i\vec{\beta}$，
将 $w_1\vec{\rho}_1$ 倍加到 $w_i\vec{\rho}_1$，等等，
从而得到 $\vec{v}$ 和 $\vec{w}$ 的展开式。
%</pf:DetsMultilinear3>
\end{proof}

多重线性使我们能把一个行列式展开成一些行列式的和，
其中每一个都涉及一个简单的矩阵。

\begin{example}
利用多重线性的性质~(1)
拆分第一行
\begin{equation*}
  \begin{vmat}[r]
     2  &1  \\
     4  &3
  \end{vmat}
  =
  \begin{vmat}[r]
     2  &0  \\
     4  &3
  \end{vmat}
  +
  \begin{vmat}[r]
     0  &1  \\
     4  &3
  \end{vmat}
\end{equation*}
然后再用 (1) 沿第二行拆分每一个。
\begin{equation*}
  =\begin{vmat}[r]
     2  &0  \\
     4  &0
  \end{vmat}
  +
  \begin{vmat}[r]
     2  &0  \\
     0  &3
  \end{vmat}
  +
  \begin{vmat}[r]
     0  &1  \\
     4  &0
  \end{vmat}
  +
  \begin{vmat}[r]
     0  &1  \\
     0  &3
  \end{vmat}
\end{equation*}
结果是四个行列式。
在这四个行列式的每一行中，
都只有来自原矩阵的一个元素。
\end{example}

\begin{example}
同样地，
一个 \( \nbyn{3} \) 行列式分离成
许多更简单行列式的和。
沿第一行拆分产生三个行列式
（我们高亮了 $1,3$ 位置的零，使它在视觉上
与作为拆分一部分出现的零区分开）。
\begin{equation*}
  \begin{vmat}[r]
     2              &1  &-1  \\
     4              &3  &\highlight{0}  \\
     2              &1  &5
  \end{vmat}
  =
  \begin{vmat}[r]
     2              &0  &0   \\
     4              &3  &\highlight{0}  \\
     2              &1  &5
  \end{vmat}
  +
  \begin{vmat}[r]
     0              &1  &0   \\
     4              &3  &\highlight{0}   \\
     2              &1  &5
  \end{vmat}
  +
  \begin{vmat}[r]
     0              &0  &-1  \\
     4              &3  &\highlight{0}  \\
     2  &1  &5
  \end{vmat}                       
\end{equation*}
反过来，上面每一个都沿第二行拆分成三个。
然后这九个每一个再沿第三行拆分成三个。
结果是二十七个行列式，
使得每一行都只包含来自初始矩阵的一个元素。
\begin{equation*}
  =
  \begin{vmat}[r]
     2              &0  &0   \\
     4              &0  &0   \\
     2              &0  &0
  \end{vmat}
  +
  \begin{vmat}[r]
     2  &0  &0   \\
     4  &0  &0   \\
     0  &1  &0
  \end{vmat}
  +
  \begin{vmat}[r]
     2  &0  &0   \\
     4  &0  &0   \\
     0  &0  &5
  \end{vmat}
  +
  \begin{vmat}[r]
     2  &0  &0   \\
     0  &3  &0   \\
     2  &0  &0
  \end{vmat}
  +\dots+
  \begin{vmat}[r]
     0  &0  &-1  \\
     0  &0  &\highlight{0}  \\
     0  &0  &5
  \end{vmat}
\end{equation*}
\end{example}

于是多重线性会将一个 \( \nbyn{n} \) 行列式展开成
\( n^n \) 个行列式的和，其中
每个行列式的每一行都只包含来自
初始矩阵的一个元素。

在这个展开中，尽管项很多，
但大多数的行列式值为零。

\begin{example} \label{ex:SamplePermExp}
在来自前面展开式的这些例子中，
原矩阵的两个元素位于同一列。
\begin{equation*}
  \begin{vmat}[r]
     2               &0  &0   \\
     4               &0  &0  \\
     0               &1  &0
  \end{vmat}
  \qquad
  \begin{vmat}[r]
     0               &0  &-1  \\
     0               &3  &0  \\
     0               &0  &5
  \end{vmat}
  \qquad\
  \begin{vmat}[r]
     0               &1  &0   \\
     0               &0  &\highlight{0}  \\
     0               &0  &5
  \end{vmat}
\end{equation*}
例如，
在第一个矩阵中，$2$ 和 $4$ 都来自原矩阵的第一列。
在第二个矩阵中，$-1$ 和~$5$ 都来自第三列。
而在第三个矩阵中，$0$ 和~$5$ 都来自第三列。
任何这样的矩阵都是奇异的，因为
其中一行是另一行的倍数。
因此，根据 \nearbylemma{le:IdenRowsDetZero}，任何这样的行列式都为零。

有了这一观察，
上面把 \( \nbyn{3} \) 行列式展开成
二十七个行列式之和的做法就化简为这六个行列式之和：
其中来自原矩阵的元素每行一个，
且每列也只有一个。
\begin{align*}
  \begin{vmat}[r]
     2  &1  &-1  \\
     4  &3  &\highlight{0}  \\
     2  &1  &5
  \end{vmat}
  &=\begin{vmat}[r]
     2  &0  &0   \\
     0  &3  &0   \\
     0  &0  &5
  \end{vmat}
  +
   \begin{vmat}[r]
     2  &0  &0   \\
     0  &0  &\highlight{0}   \\
     0  &1  &0
  \end{vmat}                      \\
  &\quad\hbox{}+\begin{vmat}[r]
     0  &1  &0   \\
     4  &0  &0   \\
     0  &0  &5
  \end{vmat}
  +
   \begin{vmat}[r]
     0  &1  &0   \\
     0  &0  &\highlight{0}   \\
     2  &0  &0
  \end{vmat}                      \\
  &\quad\hbox{}+\begin{vmat}[r]
     0  &0  &-1  \\
     4  &0  &0   \\
     0  &1  &0
  \end{vmat}
  +
   \begin{vmat}[r]
     0  &0  &-1  \\
     0  &3  &0    \\
     2  &0  &0
  \end{vmat}                      
\end{align*}
在那个展开式中我们可以把标量提出来。
\begin{align*}
  &=(2)(3)(5)\begin{vmat}[r]
               1  &0  &0  \\
               0  &1  &0  \\
               0  &0  &1
            \end{vmat}
  +(2)(\highlight{0})(1)\begin{vmat}[r]
               1  &0  &0  \\
               0  &0  &1  \\
               0  &1  &0
            \end{vmat}                     \\
  &\quad\hbox{}+(1)(4)(5)\begin{vmat}[r]
               0  &1  &0  \\
               1  &0  &0  \\
               0  &0  &1
            \end{vmat}
  +(1)(\highlight{0})(2)\begin{vmat}[r]
               0  &1  &0  \\
               0  &0  &1  \\
               1  &0  &0
            \end{vmat}                       \\
  &\quad\hbox{}+(-1)(4)(1)\begin{vmat}[r]
               0  &0  &1  \\
               1  &0  &0  \\
               0  &1  &0
            \end{vmat}
  +(-1)(3)(2)\begin{vmat}[r]
               0  &0  &1  \\
               0  &1  &0  \\
               1  &0  &0
            \end{vmat}                       
\end{align*}
为完成计算，通过行对换把那六个行列式
化为单位矩阵，
并记录符号的变化。
\begin{align*}
  &=30\cdot (+1)+0\cdot (-1)  \\
  &\quad\hbox{}+20\cdot (-1)+0\cdot (+1) \\
  &\quad \hbox{}-4\cdot (+1)-6\cdot (-1)=12
\end{align*}
\end{example}

这个例子体现了本小节新的计算方案。
多重线性把一个行列式展开成
许多单独的行列式，每个行列式的每一行有来自
原矩阵的一个元素。
其中大多数都
有一行是另一行的倍数，
所以我们把它们略去。
剩下的是那些每行每列各有一个原矩阵元素的行列式。
提出标量进一步把我们必须计算的行列式
化简为每行每列恰有一个元素的、所有元素都是 $1$ 的矩阵。

%<*NotationForPermutationMatrices>
回忆定义~Three.IV.\ref{df:PermutationMatrix}：
一个 \definend{置换矩阵}\index{permutation matrix}%
\index{matrix!permutation}
是方阵，其元素除每行每列恰有一个 $1$ 外都是 $0$。
现在我们引入置换矩阵的一种记号。
%</NotationForPermutationMatrices>

\begin{definition}  \label{df:permutation}
%<*df:permutation>
一个 \definend{\( n \)-置换}\index{permutation}
是前~$n$ 个正整数上的一个函数
$\map{\phi}{\set{1,\ldots,n}}{\set{1,\ldots,n}}$，
它是单射且满射的。
%</df:permutation>
\end{definition}

在一个置换中，每个数
$1$, \ldots, $n$ 都作为且仅作为一个输入的输出出现。
我们可以把置换记为一个序列
$\phi=\sequence{\phi(1),\phi(2),\ldots,\phi(n)}$。

\begin{example} \label{ex:AllTwoThreePerms}
%<*ex:AllTwoPerms>
$2$-置换是函数
$\map{\phi_1}{\set{1,2}}{\set{1,2}}$
由 $\phi_1(1)=1$，$\phi_1(2)=2$ 给出，
以及
$\map{\phi_2}{\set{1,2}}{\set{1,2}}$
由 $\phi_2(1)=2$，$\phi_2(2)=1$ 给出。
序列记号更简洁：
\( \phi_1=\sequence{1,2} \) 与 \( \phi_2=\sequence{2,1} \)。
%</ex:AllTwoPerms>
\end{example}

\begin{example} \label{ex:AllThreePerms}
%<*ex:AllThreePerms>
在序列记号下，$3$-置换是
\( \phi_1=\sequence{1,2,3} \)、
\( \phi_2=\sequence{1,3,2} \)、
\( \phi_3=\sequence{2,1,3} \)、
\( \phi_4=\sequence{2,3,1} \)、
\( \phi_5=\sequence{3,1,2} \)，以及
\( \phi_6=\sequence{3,2,1} \)。
%</ex:AllThreePerms>
\end{example}

%<*AssociatePermutationWithPermutationMatrix>
我们把除第 $j$ 个元素处为 $1$ 外全为 $0$ 的行向量记为 $\iota_j$，
于是宽度为四的 $\iota_2$
是~$\rowvec{0 &1 &0 &0}$。
现在我们的置换矩阵记号是：
对任意 $\phi=\sequence{\phi(1),\ldots,\phi(n)}$，
与之相伴的矩阵的行依次为 $\iota_{\phi(1)}$, \ldots, $\iota_{\phi(n)}$。
%</AssociatePermutationWithPermutationMatrix>
例如，与 $4$-置换 \( \phi=\sequence{3,2,1,4} \) 相伴的
是其行依次为相应 $\iota$ 的矩阵。
\begin{equation*}
   P_{\phi}=
   \begin{mat}
      \iota_{3} \\
      \iota_{2} \\
      \iota_{1} \\
      \iota_{4} 
   \end{mat}=
  \begin{mat}[r]
     0  &0  &1  &0  \\
     0  &1  &0  &0  \\
     1  &0  &0  &0  \\
     0  &0  &0  &1
  \end{mat}
\end{equation*}

\begin{example} \label{ex:TwoThreePermsAndMatrices}
这些是 \nearbyexample{ex:AllTwoThreePerms} 中列出的
$2$-置换的置换矩阵。
\begin{equation*}
  P_{\phi_1}
  =\begin{mat}
      \iota_1 \\
      \iota_2 
   \end{mat}
  =\begin{mat}[r]
    1  &0         \\
    0  &1   
  \end{mat}
  \qquad
  P_{\phi_2}
   =\begin{mat}
      \iota_2 \\
      \iota_1 
   \end{mat}
   =\begin{mat}[r]
    0  &1         \\
    1  &0   
  \end{mat}
\end{equation*}
例如，$P_{\phi_2}$ 的第一行是
$\iota_{\phi_2(1)}=\iota_2$，第二行是 $\iota_{\phi_2(2)}=\iota_1$。
\end{example}

\begin{example}
考虑 $3$-置换
\( \phi_5=\sequence{3,1,2} \)。
置换矩阵  
$P_{\phi_5}$ 的行依次为 $\iota_{\phi_5(1)}=\iota_3$、
$\iota_{\phi_5(2)}=\iota_1$ 与 $\iota_{\phi_5(3)}=\iota_2$。\begin{equation*}
  % P_{\phi_2}
  %  =\begin{mat}
  %     \iota_1 \\
  %     \iota_3 \\
  %     \iota_2 
  %  \end{mat}
  % =\begin{mat}[r]
  %    1      &0        &0        \\
  %    0      &0        &1        \\
  %    0      &1        &0
  % \end{mat}
  % \qquad
  P_{\phi_5}
   % =\begin{mat}
   %    \iota_3 \\
   %    \iota_1 \\
   %    \iota_2 
   % \end{mat}
  =\begin{mat}[r]
     0      &0        &1        \\
     1      &0        &0        \\
     0      &1        &0
  \end{mat}
\end{equation*}
\end{example}

\begin{definition} \label{df:PermutationExpansion}
%<*df:PermutationExpansion>
行列式的 \definend{置换展开}\index{determinant!permutation expansion}%
\index{permutation expansion}
是
\begin{equation*}
   \begin{vmat}
      t_{1,1}  &t_{1,2}  &\ldots  &t_{1,n}  \\
      t_{2,1}  &t_{2,2}  &\ldots  &t_{2,n}  \\
               &\vdots                      \\
      t_{n,1}  &t_{n,2}  &\ldots  &t_{n,n}
   \end{vmat}
   =
   \begin{array}[t]{l}
      t_{1,\phi_1(1)}t_{2,\phi_1(2)}\cdots
           t_{n,\phi_1(n)}\deter{P_{\phi_1}}       \\[.5ex]
      \>\hbox{}+t_{1,\phi_2(1)}t_{2,\phi_2(2)}\cdots
           t_{n,\phi_2(n)}\deter{P_{\phi_2}}       \\
      \>\vdotswithin{+}                              \\
      \>\hbox{}+t_{1,\phi_k(1)}t_{2,\phi_k(2)}\cdots
           t_{n,\phi_k(n)}\deter{P_{\phi_k}} 
   \end{array}
\end{equation*}
其中 \( \phi_1,\ldots,\phi_k \) 是所有的 \( n \)-置换。
%</df:PermutationExpansion>
\end{definition}
%<*SummationForPermutationExpansion>
我们可以把这个公式改述为
\definend{求和
记号}\index{summation notation, for permutation expansion}
\begin{equation*}
  \deter{T}=\!\!
  \sum_{\text{置换\ }\phi}\!\!\!\!
     t_{1,\phi(1)}t_{2,\phi(2)}\cdots t_{n,\phi(n)}
                                 \deter{P_{\phi}}
\end{equation*}
读作，
"对所有的置换 \( \phi \) 求和，项形如
\( t_{1,\phi(1)}t_{2,\phi(2)}\cdots t_{n,\phi(n)} \deter{P_{\phi}} \)。"
%</SummationForPermutationExpansion>


\begin{example}
熟悉的 $\nbyn{2}$ 行列式公式可由上面推出
\begin{align*}
  \begin{vmat}
    t_{1,1}  &t_{1,2} \\
    t_{2,1}  &t_{2,2}
  \end{vmat}
  &=
  t_{1,1}t_{2,2}\cdot\deter{P_{\phi_1}}
  +
  t_{1,2}t_{2,1}\cdot\deter{P_{\phi_2}}      \\[-1.5ex]     
  &=
  t_{1,1}t_{2,2}\cdot\begin{vmat}[r]
           1  &0 \\
           0  &1
         \end{vmat}
  +
  t_{1,2}t_{2,1}\cdot\begin{vmat}[r]
           0  &1 \\
           1  &0
         \end{vmat}               \\
  &=t_{1,1}t_{2,2}-t_{1,2}t_{2,1}
\end{align*}
$\nbyn{3}$ 公式也是如此。
\begin{align*}
  \begin{vmat}
    t_{1,1}  &t_{1,2}  &t_{1,3} \\
    t_{2,1}  &t_{2,2}  &t_{2,3} \\
    t_{3,1}  &t_{3,2}  &t_{3,3} 
  \end{vmat}
  &=
  \begin{aligned}[t]
    &t_{1,1}t_{2,2}t_{3,3}\deter{P_{\phi_1}}
     +t_{1,1}t_{2,3}t_{3,2}\deter{P_{\phi_2}}
     +t_{1,2}t_{2,1}t_{3,3}\deter{P_{\phi_3}} \\
    &\hspace*{0em}
     +t_{1,2}t_{2,3}t_{3,1}\deter{P_{\phi_4}}
     +t_{1,3}t_{2,1}t_{3,2}\deter{P_{\phi_5}}
     +t_{1,3}t_{2,2}t_{3,1}\deter{P_{\phi_6}}
  \end{aligned}                                      \\
  &=
  \begin{aligned}[t]
    &t_{1,1}t_{2,2}t_{3,3}
     -t_{1,1}t_{2,3}t_{3,2}
     -t_{1,2}t_{2,1}t_{3,3}  \\
    &\hspace*{0em}
     +t_{1,2}t_{2,3}t_{3,1}
     +t_{1,3}t_{2,1}t_{3,2}
     -t_{1,3}t_{2,2}t_{3,1}
  \end{aligned}
\end{align*}
\end{example}

用置换展开计算行列式通常比用高斯消元法
耗时更长。
不过，我们将用它来证明
行列式函数存在。
这个证明很长，所以我们在这里只陈述结果，
而把证明留到下一小节。

\begin{theorem} \label{th:DetsExist}
\index{determinant!exists} 
%<*th:DetsExist>
对每个 $n$，都存在一个 $\nbyn{n}$ 行列式函数。
%</th:DetsExist>
\end{theorem}

下一小节还将给出
下一个结果的证明（把它们放在一起是因为两个证明有重叠）。

\begin{theorem}  \label{th:DeterminantOfAMatrixEqualsDeterminantOfTranspose}
\index{transpose!determinant}
%<*th:DeterminantOfAMatrixEqualsDeterminantOfTranspose>
一个矩阵的行列式等于其转置的行列式。
%</th:DeterminantOfAMatrixEqualsDeterminantOfTranspose>
\end{theorem}

由于这个定理，
虽然迄今为止我们都是按行陈述行列式的结果，
但所有这些结果按列也同样成立。

\begin{corollary} \label{cor:ColSwapChgSign} \label{cor:DetsMultiInCols}
%<*co:ColSwapChgSign>
  有两列相等的矩阵是奇异的。
  列对换改变行列式的符号。
  行列式关于其列是多重线性的。
%</co:ColSwapChgSign>
\end{corollary}

\begin{proof}
%<*pf:ColSwapChgSign>
对于第一个命题，
转置该矩阵得到一个具有相同行列式的矩阵，
且它有两行相等，从而行列式为零。
其余两个命题按同样的方式证明。
%</pf:ColSwapChgSign>
\end{proof}

我们以一段总结来结束本小节：
行列式函数存在、唯一，并且我们知道如何计算它们。
至于行列式究竟是关于什么的，下面几行文字
\cite{Kemp} 或许有助于我们记住它。
\begin{verse} \small
行列式为零，         \\*
\ 解：解无穷多或无解。   \\*
行列式非零，         \\*
\ 解：恰有一解。
\end{verse}




\begin{exercises}
  \item[]\textit{这里总结我们对
    $2$-\hbox{}与 $3$-置换所用的记号。
    \begin{center}
      \begin{tabular}[t]{c|cc}
        $i$          &$1$      &$2$    \\
        \hline
        $\phi_1(i)$  &$1$      &$2$     \\
        $\phi_2(i)$  &$2$      &$1$
      \end{tabular}
      \qquad
      \begin{tabular}[t]{c|ccc}
        $i$          &$1$     &$2$   &$3$    \\
        \hline
        $\phi_1(i)$  &$1$     &$2$   &$3$    \\
        $\phi_2(i)$  &$1$     &$3$   &$2$    \\
        $\phi_3(i)$  &$2$     &$1$   &$3$    \\
        $\phi_4(i)$  &$2$     &$3$   &$1$    \\
        $\phi_5(i)$  &$3$     &$1$   &$2$    \\
        $\phi_6(i)$  &$3$     &$2$   &$1$
      \end{tabular}
    \end{center}  }
  \recommended \item 对该矩阵，求与每个 $3$-置换相联系的项。
    \begin{equation*}
      M=\begin{mat}
        1 &2 &3 \\
        4 &5 &6 \\
        7 &8 &9
      \end{mat}
    \end{equation*}
    也就是说，填写下表中的其余部分。
      \begin{center} \small
        \begin{tabular}{r|cccccc}
          置换 $\phi_i$
            &$\phi_1$ &$\phi_2$ &$\phi_3$ &$\phi_4$ &$\phi_5$ &$\phi_6$  \\
          \hline
          项 $m_{1,\phi_i(1)}m_{2,\phi_i(2)}m_{3,\phi_i(3)}$
           &$1\cdot 5\cdot 9$
           &
           &
           &
           &
           &
        \end{tabular}
      \end{center}
    \begin{answer}
      记该矩阵为~$M$。
      \begin{center}
        \begin{tabular}{r|cccccc}
          置换
            &$\phi_1$ &$\phi_2$ &$\phi_3$ &$\phi_4$ &$\phi_5$ &$\phi_6$  \\
          \hline
          项
           &$1\cdot 5\cdot 9$
           &$1\cdot 6\cdot 8$
           &$2\cdot 4\cdot 9$
           &$2\cdot 6\cdot 7$
           &$3\cdot 4\cdot 8$
           &$3\cdot 5\cdot 7$
        \end{tabular}
      \end{center}
    \end{answer}
  \recommended \item 对每个 $3$-置换~$\phi$，求 $|P_\phi|$。
    \begin{answer}
      我们可以把每个 $P_{\phi}$ 经过行对换变为单位矩阵。
      如果对换的次数是偶数，那么它的行列式为~$+1$；
      而如果对换的次数是奇数，那么行列式为~$-1$。
      \begin{center}
        \begin{tabular}{r|cccccc}
          置换
            &$\phi_1$ &$\phi_2$ &$\phi_3$ &$\phi_4$ &$\phi_5$ &$\phi_6$  \\
          \hline
          对换次数 &$0$
                  &$1$
                  &$1$
                  &$2$
                  &$2$
                  &$1$  \\
          $|P_{\phi_i}|$
           &$+1$
           &$-1$
           &$-1$
           &$+1$
           &$+1$
           &$-1$
        \end{tabular}
      \end{center}
      \textit{注。}
      在上表中，
      我们只是不断地对换，直到找到一个把我们带回到`$1,2,3$'的数为止。
      但不同的对换次数是否可能？
      如果一个人找到 $2$ 次对换而另一个人找到
      $4$ 次，这没有问题，因为两者给出的行列式都是~$+1$。
      但如果我们能用两种不同的方式对换，其中一种次数为偶数而另一种为奇数，
      那就有问题了。
      下一节将表明偶数与奇数混合的情形不可能出现。
    \end{answer}
  \item 由 $\nbyn{2}$ 公式知该行列式为~$7$。
    请用置换展开计算它。
    \begin{equation*}
      \begin{vmatrix}
        2 &3 \\
        1 &5
      \end{vmatrix}
    \end{equation*}
    \begin{answer}
      \begin{align*}
        \begin{vmat}
          2  &3 \\
          1  &5
        \end{vmat}
        &=
        2\cdot 5\cdot\deter{P_{\phi_1}}
        +
        1\cdot 3\cdot\deter{P_{\phi_2}}      \\
        &=
        2\cdot 5\cdot\begin{vmat}[r]
                           1  &0 \\
                           0  &1
                         \end{vmat}
        +
        1\cdot 3\cdot\begin{vmat}[r]
                           0  &1 \\
                           1  &0
                         \end{vmat}               \\
        &=2\cdot 5\cdot 1+1\cdot 3\cdot (-1)=7
        \end{align*}
    \end{answer}
  \item 该行列式为 $0$，因为前两行之和等于第三行。
    请用置换展开计算该行列式。
    \begin{equation*}
      \begin{vmatrix}
        -1 &0 &1 \\
        3  &1 &4 \\
        2  &1 &5
      \end{vmatrix}
    \end{equation*}
    \begin{answer}
          \begin{align*}
            \begin{vmat}[r]
              -1  &0  &1  \\
               3  &1  &4  \\
               2  &1  &5
            \end{vmat}
            &=
            \begin{aligned}[t]
              &(-1)(1)(5)\deter{P_{\phi_1}}
              +(-1)(4)(1)\deter{P_{\phi_2}}
              +(0)(3)(5)\deter{P_{\phi_3}}  \\
              &\hbox{}\quad\hbox{}
              +(0)(4)(2)\deter{P_{\phi_4}}
              +(1)(3)(1)\deter{P_{\phi_5}}
              +(1)(1)(2)\deter{P_{\phi_6}}
            \end{aligned}                        \\
            &=
            \begin{aligned}[t]
              &(-1)(1)(5)\cdot 1
              +(-1)(4)(1)\cdot (-1)
              +(0)(3)(5)\cdot (-1)  \\
              &\hbox{}\quad\hbox{}
              +(0)(4)(2)\cdot 1
              +(1)(3)(1)\cdot 1
              +(1)(1)(2)\cdot (-1)
            \end{aligned}                        \\
            &=-5+4+0+0+3-2=0
          \end{align*}
    \end{answer}
  \recommended \item
    用置换展开计算行列式。
    \begin{exparts*}
      \partsitem
        $\begin{vmat}[r]
          1  &2  &3  \\
          4  &5  &6  \\
          7  &8  &9
        \end{vmat}$
      \partsitem
        $\begin{vmat}[r]
          2  &2  &1  \\
          3  &-1 &0  \\
          -2 &0  &5
        \end{vmat}$
    \end{exparts*}
    \begin{answer}
      \begin{exparts}
        \partsitem 该矩阵是奇异的。
          \begin{align*}
            \begin{vmat}[r]
              1  &2  &3  \\
              4  &5  &6  \\
              7  &8  &9
            \end{vmat}
            &=
            \begin{aligned}[t]
              &(1)(5)(9)\deter{P_{\phi_1}}
              +(1)(6)(8)\deter{P_{\phi_2}}
              +(2)(4)(9)\deter{P_{\phi_3}}  \\
              &\hbox{}\quad\hbox{}
              +(2)(6)(7)\deter{P_{\phi_4}}
              +(3)(4)(8)\deter{P_{\phi_5}}
              +(7)(5)(3)\deter{P_{\phi_6}}
            \end{aligned}                        \\
            &=0
          \end{align*}
        \partsitem 该矩阵是非奇异的。
          \begin{align*}
            \begin{vmat}[r]
              2  &2  &1  \\
              3  &-1 &0  \\
              -2 &0  &5
            \end{vmat}
            &=
            \begin{aligned}[t]
              &(2)(-1)(5)\deter{P_{\phi_1}}
              +(2)(0)(0)\deter{P_{\phi_2}}
              +(2)(3)(5)\deter{P_{\phi_3}} \\
              &\hbox{}\quad\hbox{}
              +(2)(0)(-2)\deter{P_{\phi_4}}
              +(1)(3)(0)\deter{P_{\phi_5}}
              +(-2)(-1)(1)\deter{P_{\phi_6}}
            \end{aligned}                      \\
            &=-42
        \end{align*}
      \end{exparts}
    \end{answer}
  \recommended \item
    分别用高斯消元法与置换展开公式计算下面各行列式。
    \begin{exparts*}
      \partsitem \( \begin{vmat}[r]
                 2  &1  \\
                 3  &1
               \end{vmat}   \)
      \partsitem \( \begin{vmat}[r]
                  0  &1  &4  \\
                  0  &2  &3  \\
                  1  &5  &1
               \end{vmat}  \)
    \end{exparts*}
    \begin{answer}
      \begin{exparts}
        \partsitem 高斯消元法给出
          \begin{equation*}
             \begin{vmat}[r]
                 2  &1  \\
                 3  &1
             \end{vmat}
             =
             \begin{vmat}[r]
                 2  &1  \\
                 0  &-1/2
             \end{vmat}
             =-1
          \end{equation*}
          而置换展开给出
          \begin{equation*}
             \begin{vmat}[r]
                 2  &1  \\
                 3  &1
             \end{vmat}
             =
             \begin{vmat}[r]
                 2  &0  \\
                 0  &1
             \end{vmat} +
             \begin{vmat}[r]
                 0  &1  \\
                 3  &0
             \end{vmat}
             =
             (2)(1)\begin{vmat}[r]
                 1  &0  \\
                 0  &1
             \end{vmat} +
             (1)(3)\begin{vmat}[r]
                 0  &1  \\
                 1  &0
             \end{vmat}
             =
             -1
          \end{equation*}
        \partsitem 高斯消元法给出
          \begin{equation*}
            \begin{vmat}[r]
              0  &1  &4  \\
              0  &2  &3  \\
              1  &5  &1
            \end{vmat}
            =
            -\begin{vmat}[r]
              1  &5  &1  \\
              0  &2  &3  \\
              0  &1  &4
            \end{vmat}
            =
            -\begin{vmat}[r]
              1  &5  &1  \\
              0  &2  &3  \\
              0  &0  &5/2
            \end{vmat}
            =-5
          \end{equation*}
          而置换展开给出
          \begin{align*}
            \begin{vmat}[r]
              0  &1  &4  \\
              0  &2  &3  \\
              1  &5  &1
            \end{vmat}
            &=
            \begin{aligned}[t]
            &(0)(2)(1)\deter{P_{\phi_1}}
            +(0)(3)(5)\deter{P_{\phi_2}}
            +(1)(0)(1)\deter{P_{\phi_3}}  \\
            &\hbox{}\quad\hbox{}
            +(1)(3)(1)\deter{P_{\phi_4}}
            +(4)(0)(5)\deter{P_{\phi_5}}
            +(4)(2)(1)\deter{P_{\phi_6}}
            \end{aligned}                         \\
            &=-5
          \end{align*}
      \end{exparts}
    \end{answer}
  \recommended \item
    用置换展开公式推导 \( \nbyn{3} \) 行列式的公式。
    \begin{answer}
        仿照 \nearbyexample{ex:SamplePermExp} 可得
          \begin{align*}
             \begin{vmat}
                t_{1,1}  &t_{1,2}  &t_{1,3}  \\
                t_{2,1}  &t_{2,2}  &t_{2,3}  \\
                t_{3,1}  &t_{3,2}  &t_{3,3}
             \end{vmat}
             &=\begin{aligned}[t]
                 &t_{1,1}t_{2,2}t_{3,3}\deter{P_{\phi_1}}
                   +t_{1,1}t_{2,3}t_{3,2}\deter{P_{\phi_2}} \\
                 &\hbox{}\quad\hbox{}
                  +t_{1,2}t_{2,1}t_{3,3}\deter{P_{\phi_3}}
                             +t_{1,2}t_{2,3}t_{3,1}\deter{P_{\phi_4}} \\
                 &\hbox{}\quad\hbox{}
                  +t_{1,3}t_{2,1}t_{3,2}\deter{P_{\phi_5}}
                             +t_{1,3}t_{2,2}t_{3,1}\deter{P_{\phi_6}}
               \end{aligned}                                               \\
             &=\begin{aligned}[t]
                 & t_{1,1}t_{2,2}t_{3,3}(+1)
                   +t_{1,1}t_{2,3}t_{3,2}(-1)  \\
                 &\hbox{}\quad\hbox{}
                   +t_{1,2}t_{2,1}t_{3,3}(-1)
                    +t_{1,2}t_{2,3}t_{3,1}(+1) \\
                 &\hbox{}\quad\hbox{}
                   +t_{1,3}t_{2,1}t_{3,2}(+1)
                         +t_{1,3}t_{2,2}t_{3,1}(-1)
               \end{aligned}
          \end{align*}
     \end{answer}
  \item
    列出全部 $4$-置换。
    \begin{answer}
      下面是所有满足 $\phi(1)=1$ 的置换
      \begin{align*}
        &\phi_1=\sequence{1,2,3,4}
        \quad
        \phi_2=\sequence{1,2,4,3}
        \quad
        \phi_3=\sequence{1,3,2,4}  \\
        &\quad
        \phi_4=\sequence{1,3,4,2}
        \quad
        \phi_5=\sequence{1,4,2,3}
        \quad
        \phi_6=\sequence{1,4,3,2}
      \end{align*}
      下面是满足 $\phi(1)=1$ 的那些
      \begin{align*}
        &\phi_7=\sequence{2,1,3,4}
        \quad
        \phi_8=\sequence{2,1,4,3}
        \quad
        \phi_9=\sequence{2,3,1,4}    \\
        &\quad
        \phi_{10}=\sequence{2,3,4,1}
        \quad
        \phi_{11}=\sequence{2,4,1,3}
        \quad
        \phi_{12}=\sequence{2,4,3,1}
      \end{align*}
      下面是满足 $\phi(1)=3$ 的那些
      \begin{align*}
        &\phi_{13}=\sequence{3,1,2,4}
        \quad
        \phi_{14}=\sequence{3,1,4,2}
        \quad
        \phi_{15}=\sequence{3,2,1,4}  \\
        &\quad
        \phi_{16}=\sequence{3,2,4,1}
        \quad
        \phi_{17}=\sequence{3,4,1,2}
        \quad
        \phi_{18}=\sequence{3,4,2,1}
      \end{align*}
      以及满足 $\phi(1)=4$ 的那些。
      \begin{align*}
        &\phi_{19}=\sequence{4,1,2,3}
        \quad
        \phi_{20}=\sequence{4,1,3,2}
        \quad
        \phi_{21}=\sequence{4,2,1,3}  \\
        &\quad
        \phi_{22}=\sequence{4,2,3,1}
        \quad
        \phi_{23}=\sequence{4,3,1,2}
        \quad
        \phi_{24}=\sequence{4,3,2,1}
      \end{align*}
    \end{answer}
  \item
    置换可以看作从集合
    $\set{1,..,n}$ 到其自身的一一对应（既是单射又是满射）的函数。
    因此，每个置换都有逆。
    \begin{exparts}
      \partsitem 求每个 $2$-置换的逆。
      \partsitem 求每个 $3$-置换的逆。
    \end{exparts}
    \begin{answer}
      下列每一个都容易验证。
      \begin{exparts}
        \partsitem
          \begin{tabular}[t]{r|cc}
            置换 &$\phi_1$  &$\phi_2$ \\
             \hline
            逆     &$\phi_1$  &$\phi_2$
          \end{tabular}
        \partsitem
          \begin{tabular}[t]{r|cccccc}
            置换
              &$\phi_1$ &$\phi_2$ &$\phi_3$ &$\phi_4$ &$\phi_5$ &$\phi_6$ \\
            \hline
            逆
              &$\phi_1$ &$\phi_2$ &$\phi_3$ &$\phi_5$ &$\phi_4$ &$\phi_6$
          \end{tabular}
      \end{exparts}
    \end{answer}
  \item
     证明：\( f \) 是多重线性的当且仅当对所有
     \( \vec{v},\vec{w}\in V \) 与 \( k_1,k_2\in\Re \)，下式成立。
      \begin{equation*}
         f(\vec{\rho}_1,\dots,k_1\vec{v}_1+k_2\vec{v}_2,
           \dots,\vec{\rho}_n)
         =
         k_1f(\vec{\rho}_1,\dots,\vec{v}_1,\dots,\vec{\rho}_n)+
         k_2f(\vec{\rho}_1,\dots,\vec{v}_2,\dots,\vec{\rho}_n)
       \end{equation*}
       \begin{answer}
         对于「如果」部分，取 $k_1=k_2=1$ 即得
         \nearbydefinition{def:multilinear} 的第一个条件，
         取 $k_2=0$ 即得第二个条件。

         「仅当」部分也是常规推导。
         由
         $
           f(\vec{\rho}_1,\dots,k_1\vec{v}_1+k_2\vec{v}_2,
             \dots,\vec{\rho}_n)
         $
         出发，\nearbydefinition{def:multilinear} 的第一个条件给出
         $
           =
           f(\vec{\rho}_1,\dots,k_1\vec{v}_1,\dots,\vec{\rho}_n)+
           f(\vec{\rho}_1,\dots,k_2\vec{v}_2,\dots,\vec{\rho}_n)
         $
         再对第二个条件应用两次，即得结论。
       \end{answer}
  \item
    如果把定义中的性质~(4) 改成 \( \deter{I}=2 \)，行列式会如何变化？
    \begin{answer}
       它们都会变成原来的两倍。
    \end{answer}
  \item
    验证 \nearbycorollary{cor:ColSwapChgSign} 中的第二、第三个论断。
    \begin{answer}
      对于第二个论断，
      给定一个矩阵，先转置，再对换行，然后转置回来。
      所得结果是对换了列，而行列式改变了因子
      \( -1 \)。
      第三个论断类似：~给定
      一个矩阵，先转置，对如今已是行的东西应用多重线性，
      然后把所得的各个矩阵转置回来。
    \end{answer}
  \recommended \item
    证明：若 \( \nbyn{n} \) 矩阵的行列式非零，
    则我们能将任一列向量
    \( \vec{v}\in\Re^n \) 表示为该矩阵各列的线性组合。
    \begin{answer}
      行列式非零的 \( \nbyn{n} \) 矩阵的秩为
      \( n \)，所以它的各列构成 \( \Re^n \) 的一组基。
    \end{answer}
  \item
    \cite{Strang}
    判断真假：元素只取 $0$ 或 $1$ 的矩阵，其行列式必等于零、一或负一。
    \begin{answer}
      假。
      \begin{equation*}
        \begin{vmat}[r]
          0 &1  &1  \\
          1 &0  &1 \\
          1  &1 &0
        \end{vmat}
        =2
      \end{equation*}
     \end{answer}
  \item
    \begin{exparts}
      \partsitem 证明 \( \nbyn{5} \) 矩阵的置换
         展开公式中有 $120$ 项。
      \partsitem
         若 \( 1,2 \) 元为零，确定必为零的项有多少个？
    \end{exparts}
    \begin{answer}
      \begin{exparts}
        \partsitem 第一行元素的列下标有
          五种选择。
          于是第二行元素的列下标有
          四种选择。
          依此类推，我们得到 $5\cdot 4\cdot 3\cdot 2\cdot 1=120$。
          （另见下一题。）
        \partsitem 一旦选定了第一行的第二列，
          其余元素可以有 \( 4\cdot 3\cdot 2\cdot 1=24 \)
          种选法。
      \end{exparts}
    \end{answer}
  \item
    共有多少个 \( n \)-置换？
    \begin{answer}
       \( n\cdot(n-1)\cdots 2\cdot 1=n! \)
    \end{answer}
  \item 证明置换矩阵的逆是它的转置。
    \begin{answer}
      \cite{SchmidtSO}
      我们将证明 $P\trans{P}=I$；$\trans{P}P=I$ 的论证类似。
      $P\trans{P}$ 的 $i,j$ 元是形如
      $p_{i,k}q_{k,j}$ 的项之和，其中 $\trans{P}$ 的元素用 $q$ 记，
      即 $q_{k,j}=p_{j,k}$。
      于是 $P\trans{P}$ 的 $i,j$ 元是和
      \(
        \sum_{k=1}^n p_{i,k}p_{j,k}
      \)。
      但 $p_{i,k}$ 通常为 $0$，所以 $P_{i,k}P_{j,k}$ 通常为 $0$。
      $P_{i,k}$ 非零仅当它为 $1$，
      而此时再无其他 $i^\prime\neq i$ 使 $P_{i^\prime,k}$
      非零（$i$ 是第~$k$ 列中含 $1$ 的唯一一行）。
      换句话说，
      \begin{equation*}
        \sum_{k=1}^n p_{i,k}p_{j,k}=
        \begin{cases}
          1  &i=j  \\
          0  &\text{其他}
        \end{cases}
      \end{equation*}
      而这正是单位矩阵各元素的公式。
    \end{answer}
  \item
    矩阵 \( A \) 称为
    \definend{反对称的}\index{matrix!skew-symmetric}%
    \index{skew-symmetric}，若 \( \trans{A}=-A \)，
    如下矩阵。
    \begin{equation*}
      A=\begin{mat}[r]
          0  &3  \\
         -3  &0
        \end{mat}
    \end{equation*}
    证明：行列式非零的 \( \nbyn{n} \) 反对称矩阵
    仅在 \( n \) 为偶数时存在。
    \begin{answer}
      在
      \( \deter{A}=\deter{\trans{A}}=\deter{-A}=(-1)^n\deter{A} \)
      中指数 $n$ 必须为偶数。
    \end{answer}
  \recommended \item
    为保证 \( \nbyn{4} \) 矩阵的行列式为零，所需零元素的最少个数是多少？
    这些零元素应如何放置？
    \begin{answer}
      证明三个零的任何放置都不够是常规性的。
      四个零确实够了；把它们全放在同一
      行或同一列。
    \end{answer}
  \item
    设有 \( n \) 个数据点
    \( (x_1,y_1),(x_2,y_2),\dots\,,(x_n,y_n) \)，要找一条经过这些点的
    多项式 \( p(x)=a_{n-1}x^{n-1}+a_{n-2}x^{n-2}+\dots+a_1x+a_0 \)，
    那么我们可以代入这些点，得到一个 \( n \)~个方程/\( n \)~个未知数的
    线性方程组。
    该方程组的系数矩阵是
    \definend{Vandermonde 矩阵}。\index{matrix!Vandermonde}%
    \index{Vandermonde matrix}
    证明该系数矩阵的转置的行列式
    \begin{equation*}
      \begin{vmat}
         1       &1       &\ldots   &1       \\
         x_1     &x_2     &\ldots   &x_n     \\
         {x_1}^2 &{x_2}^2 &\ldots   &{x_n}^2 \\
                 &\vdots                     \\
         {x_1}^{n-1} &{x_2}^{n-1}   &\ldots   &{x_n}^{n-1}
      \end{vmat}
    \end{equation*}
    等于对所有满足 \( i<j \) 的下标 \( i,j\in\set{1,\dots,n} \)，
    形如 \( x_j-x_i \) 的项的乘积。
    （这表明，
    该行列式为零、线性方程组无解，
    当且仅当数据中的 \( x_i \) 不两两不同。）
    \begin{answer}
      $n=3$ 的情形表明了该怎么做。
      行倍加变换
      $-x_1\rho_2+\rho_3$ 与 $-x_1\rho_1+\rho_2$
      给出
      \begin{multline*}
        \begin{vmat}
          1     &1     &1     \\
          x_1   &x_2   &x_3   \\
          x_1^2 &x_2^2 &x_3^2
        \end{vmat}
        =
        \begin{vmat}
          1     &1             &1             \\
          x_1   &x_2           &x_3           \\
          0     &(-x_1+x_2)x_2 &(-x_1+x_3)x_3
        \end{vmat}                                           \\
        =
        \begin{vmat}
          1     &1             &1             \\
          0     &-x_1+x_2      &-x_1+x_3      \\
          0     &(-x_1+x_2)x_2 &(-x_1+x_3)x_3
        \end{vmat}
      \end{multline*}
      再作行倍加变换 $x_2\rho_2+\rho_3$ 便得
      所要的结果。
      \begin{equation*}
        =
        \begin{vmat}
          1     &1             &1                     \\
          0     &-x_1+x_2      &-x_1+x_3              \\
          0     &0             &(-x_1+x_3)(-x_2+x_3)
        \end{vmat}
        =(x_2-x_1)(x_3-x_1)(x_3-x_2)
      \end{equation*}
    \end{answer}
  \item
    我们可以把矩阵分成
    \definend{块}，\index{block matrix}\index{matrix!block} %
    如此处，
    \begin{equation*}
      \begin{pmat}{rr|r}
          1  &2   &0  \\
          3  &4   &0  \\  \cline{1-3}
          0  &0   &-2
       \end{pmat}
    \end{equation*}
    其中显示出四个块：左上与右下的方形的 $\nbyn{2}$ 与 $\nbyn{1}$ 块，
    以及右上与左下的零块。
    证明：若一个矩阵可以分块为
    \begin{equation*}
      T=
      \begin{pmat}{c|c}
          J   &Z_2  \\  \cline{1-2}
          Z_1 &K
       \end{pmat}
    \end{equation*}
    其中 $J$ 与 $K$ 是方阵，而 $Z_1$ 与 $Z_2$ 全为零，
    则 \( \deter{T}=\deter{J}\cdot\deter{K} \)。
    \begin{answer}
      设 \( T \) 为 \( \nbyn{n} \) 的，
      \( J \) 为 \( \nbyn{p} \) 的，
      \( K \) 为 \( \nbyn{q} \) 的。
      应用置换展开公式
      \begin{equation*}
        \deter{T}=
        \sum_{\text{\scriptsize 置换 }\phi}
                t_{1,\phi(1)}t_{2,\phi(2)}\dots t_{n,\phi(n)}
                \deter{P_{\phi}}
      \end{equation*}
      由于 \( T \) 的右上部分全为零，如果一个
      \( \phi \) 的前
      \( p \) 个列编号 \( \phi(1),\dots,\phi(p) \) 中至少有一个取自 \( p+1,\dots,n \)，
      那么由 \( \phi \) 产生的项为 \( 0 \)
      （例如，若 \( \phi(1)=n \) 则
      \( t_{1,\phi(1)}t_{2,\phi(2)}\dots t_{n,\phi(n)} \)
      为 \( 0 \)）。
      所以上式化简为对具有两段的全部置换求和：
      先重排 \( 1,\dots,p \)，随后是
      \( p+1,\dots,p+q \)。
      为看出这给出 \( \deter{J}\cdot\deter{K}  \)，将其展开。
      \begin{multline*}
         \bigg[\sum_{\substack{\text{\scriptsize 置换 }\phi_1 \\
                              \text{\scriptsize 作用于 } 1,\dots,p}}
               t_{1,\phi_1(1)}\cdots t_{p,\phi_1(p)}
               \deter{P_{\phi_1}} \bigg]                                  \\
         \cdot
         \bigg[\sum_{\substack{\text{\scriptsize 置换 }\phi_2 \\
                              \text{\scriptsize 作用于 } p+1,\dots,p+q}}
               t_{p+1,\phi_2(p+1)}\cdots t_{p+q,\phi_2(p+q)}
               \deter{P_{\phi_2}}                          \bigg]
      \end{multline*}
    \end{answer}
  \item
    证明：对任一 \( \nbyn{n} \) 矩阵 \( T \)，使矩阵 \( T-rI \) 的
    行列式为零的相异实数 \( r \) 至多有
    \( n \) 个
    （我们将在第五章用到这一结果）。
    \begin{answer}
      $n=3$ 的情形表明了会发生什么。
      \begin{equation*}
        \deter{T-rI}
        =
        \begin{vmat}
          t_{1,1}-x  &t_{1,2}   &t_{1,3}  \\
          t_{2,1}    &t_{2,2}-x &t_{2,3}  \\
          t_{3,1}    &t_{3,2}   &t_{3,3}-x
        \end{vmat}
      \end{equation*}
      置换展开中的每一项都有三个取自
      矩阵元素的因子（例如，$(t_{1,1}-x)(t_{2,2}-x)(t_{3,3}-x)$
      与 $(t_{1,1}-x)(t_{2,3})(t_{3,2})$），所以该行列式
      可以表示为 $x$ 的 $3$ 次多项式。
      这样的多项式至多有 $3$ 个根。

      一般地，置换展开表明
      行列式是一些项的和，每一项
      有 \( n \) 个因子，给出一个 $n$ 次多项式。
      \( n \) 次多项式至多有 \( n \) 个根。
    \end{answer}
  \puzzle \item
    \cite{MathMag63Q307}
    九个正的数字可以按 \( 9! \) 种方式排成 \( \nbyn{3} \) 阵列。
    求这些阵列的行列式之和。
    \begin{answer}
      \answerasgiven
      当行列式的两行对换时，行列式的
      符号改变。
      当一个三阶行列式的行作置换时，得到 \( 3 \)
      个正的与 \( 3 \) 个绝对值相等的
      负行列式。
      于是 \( 9! \) 个行列式分成 \( 9!/6 \) 组，每一组
      之和为零。
    \end{answer}
  \item
    \cite{MathMag63Q237}
    证明
    \begin{equation*}
      \begin{vmat}
        x-2  &x-3  &x-4  \\
        x+1  &x-1  &x-3  \\
        x-4  &x-7  &x-10
      \end{vmat}=0.
    \end{equation*}
    \begin{answer}
      \answerasgiven
      当任一列的元素分别从另外两列的元素中减去后，所得
      行列式的两列元素成比例，所以行列式为零。
      也就是，
      \begin{equation*}
        \begin{vmat}
          2  &1    &x-4  \\
          4  &2    &x-3  \\
          6  &3    &x-10
        \end{vmat}=
        \begin{vmat}
          1    &x-3  &-1   \\
          2    &x-1  &-2   \\
          3    &x-7  &-3
        \end{vmat}=
        \begin{vmat}
          x-2  &-1   &-2   \\
          x+1  &-2   &-4   \\
          x-4  &-3   &-6
        \end{vmat}=0.
      \end{equation*}
    \end{answer}
  \puzzle \item
    \cite{Monthly49p33}
    设 \( S \) 为一个三阶幻方的整数元素之和，\( D \) 为把该方阵视为
    行列式时的值。
    证明 \( D/S \) 是整数。
    \begin{answer}
      \answerasgiven
      设
      \begin{equation*}
        \begin{array}{ccc}
          a  &b  &c  \\
          d  &e  &f  \\
          g  &h  &i
        \end{array}
      \end{equation*}
      其幻和为 \( N=S/3 \)。
      那么
      \begin{align*}
        N
        &=(a+e+i)+(d+e+f)+(g+e+c)   \\
        &\quad\text{}-(a+d+g)-(c+f+i)=3e
      \end{align*}
      且 \( S=9e \)。
      因此，把行与列相加，
      \begin{equation*}
        D=
        \begin{vmat}
          a  &b  &c  \\
          d  &e  &f  \\
          g  &h  &i
        \end{vmat}
        =\begin{vmat}
          a  &b  &c  \\
          d  &e  &f  \\
         3e  &3e &3e
        \end{vmat}
        =\begin{vmat}
          a  &b  &3e \\
          d  &e  &3e \\
         3e  &3e &9e
        \end{vmat}
        =\begin{vmat}
          a  &b  &e  \\
          d  &e  &e  \\
          1  &1  &1
        \end{vmat}S.
      \end{equation*}
    \end{answer}
  \puzzle \item
    \cite{Monthly31p355}
    证明 Pascal 三角形左上角 \( n^2 \) 个元素构成的行列式
    \begin{equation*}
      \begin{array}{cccccc}
        1  &1  &1  &1  &.  &.  \\
        1  &2  &3  &.  &.      \\
        1  &3  &.  &.  &   &   \\
        1  &.  &.  &   &   &   \\
        .                      \\
        .
      \end{array}
    \end{equation*}
    的值为 $1$。
    \begin{answer}
      \answerasgiven
      记所论行列式为 \( D_n \)，记
      \( i \) 行 \( j \) 列的元素为 \( a_{i,j} \)。
      那么由元素的构造法则我们有
      \begin{equation*}
        a_{i,j}=a_{i,j-1}+a_{i-1,j},
        \qquad a_{1,j}=a_{i,1}=1.
      \end{equation*}
      从最后一对行开始，将 \( D_n \) 的每一行从其后一行中减去。
      经过 \( n-1 \) 次这样的减法，上面的等式表明元素
      \( a_{i,j} \) 被元素 \( a_{i,j-1} \) 所取代，而第一列中除
      \( a_{1,1}=1 \) 外的所有元素都变成零。
      现在从最后一对列开始，将每一列从其后一列中减去。
      经过这一过程，如上面的关系式所示，元素 \( a_{i,j-1} \) 被
      \( a_{i-1,j-1} \) 所取代。
      这两个运算的结果是把 \( a_{i,j} \) 换成
      \( a_{i-1,j-1} \)，并使第一行与
      第一列中的每个元素都化为零。
      于是 \( D_n=D_{n+i} \)，从而
      \begin{equation*}
        D_n=D_{n-1}=D_{n-2}=\dots=D_2=1.
      \end{equation*}
    \end{answer}
\end{exercises}














\subsectionoptional{行列式的存在性}
\textit{本小节给出前一小节中
两个结论的证明。
本小节为选读内容。
这里展开的材料只会在
若尔当标准形小节中用到，而该小节也是选读的。}

我们想要证明：对任意阶数~$n$，
$\nbyn{n}$矩阵上的行列式函数是良定义的。
前一小节给出了
置换展开公式。\index{determinant!permutation expansion}%
\index{permutation expansion}
\begin{align*}
   %\deter{T}
   %=
   \begin{vmat}
      t_{1,1}  &t_{1,2}  &\ldots  &t_{1,n}  \\
      t_{2,1}  &t_{2,2}  &\ldots  &t_{2,n}  \\
               &\vdots                      \\
      t_{n,1}  &t_{n,2}  &\ldots  &t_{n,n}
   \end{vmat}
   &=
   \begin{array}[t]{l}
      t_{1,\phi_1(1)}t_{2,\phi_1(2)}\cdots
           t_{n,\phi_1(n)}\deter{P_{\phi_1}}       \\[.75ex]
      \>\hbox{}+t_{1,\phi_2(1)}t_{2,\phi_2(2)}\cdots
           t_{n,\phi_2(n)}\deter{P_{\phi_2}}       \\[.75ex]
      \>\alignedvdots                              \\
      \>\hbox{}+t_{1,\phi_k(1)}t_{2,\phi_k(2)}\cdots
           t_{n,\phi_k(n)}\deter{P_{\phi_k}} 
   \end{array}                                                 \\[1ex]
   &=\!\!\sum_{\text{permutations\ }\phi}\!\!\!\!
     t_{1,\phi(1)}t_{2,\phi(2)}\cdots t_{n,\phi(n)}
                                 \deter{P_{\phi}}
\end{align*}
这就把证明行列式是良定义的这一
问题化归为
只需证明行列式在置换矩阵的集合上
是良定义的。

置换矩阵可以通过行对换变成单位矩阵。
因此我们可以用记录对换次数的方法
来计算它的行列式。
然而，
我们仍必须证明该结果是良定义的。
回想一下困难所在：~矩阵
\begin{equation*}
   P_{\phi}=
   \begin{mat}[r]
      0  &1  &0  &0 \\
      1  &0  &0  &0 \\
      0  &0  &1  &0 \\
      0  &0  &0  &1
   \end{mat}
\end{equation*}
的行列式既可以用一次对换
\begin{equation*}
   P_{\phi}
   \grstep{\rho_1\leftrightarrow\rho_2}
   \begin{mat}[r]
      1  &0  &0  &0 \\
      0  &1  &0  &0 \\
      0  &0  &1  &0 \\
      0  &0  &0  &1
   \end{mat}
\end{equation*}
也可以用三次对换来计算。
\begin{equation*}
   P_{\phi}
   \grstep{\rho_3\leftrightarrow\rho_1}
   % \begin{mat}[r]
   %    0  &0  &1  &0 \\
   %    1  &0  &0  &0 \\
   %    0  &1  &0  &0 \\
   %    0  &0  &0  &1
   % \end{mat}
   \repeatedgrstep{\rho_2\leftrightarrow\rho_3}
   % \begin{mat}[r]
   %    0  &0  &1  &0 \\
   %    0  &1  &0  &0 \\
   %    1  &0  &0  &0 \\
   %    0  &0  &0  &1
   % \end{mat}
   \repeatedgrstep{\rho_1\leftrightarrow\rho_3}
   \begin{mat}[r]
      1  &0  &0  &0 \\
      0  &1  &0  &0 \\
      0  &0  &1  &0 \\
      0  &0  &0  &1
   \end{mat}
\end{equation*}
这两种化简所用的对换次数都是奇数，所以这里我们认为
\( \deter{P_{\phi}}=-1 \)；
但如果有某种方法用偶数次对换来实现，
那么同一个输入就会使行列式给出
两个不同的输出。
下面，\nearbycorollary{cor:ParityInversEqParitySwaps} 将证明
这种情况不会发生\Dash
不存在这样的置换矩阵，它能以两种方式经行对换变成单位矩阵，
其中一种用偶数次对换，另一种用奇数次对换。

% So the critical step will be a way to calculate whether the 
% number of swaps that it takes could be even or odd.

\begin{definition} \label{df:Inversion}
%<*df:Inversion>
在置换 $\phi=\sequence{\ldots,k,\ldots,j,\ldots}$ 中，
满足 $k>j$ 的元素相对于其自然顺序
构成一个\definend{逆序}。
类似地，在置换矩阵中，两行
\begin{equation*}
  P_{\phi}=
  \begin{mat}
    \vdots          \\
    \iota_{k} \\
    \vdots          \\
    \iota_{j} \\
    \vdots
  \end{mat}
\end{equation*}
若满足 \( k>j \)，则处于一个\definend{逆序}中。\index{inversion}%
\index{permutation!inversions}
%</df:Inversion>
\end{definition}

\begin{example}
这个置换矩阵
\begin{equation*}
  \begin{mat}[r]
    1  &0  &0  &0 \\
    0  &0  &1  &0 \\
    0  &1  &0  &0 \\
    0  &0  &0  &1
  \end{mat}
  =
  \begin{mat}
    \iota_1  \\
    \iota_3  \\
    \iota_2  \\
    \iota_4
  \end{mat}
\end{equation*}
只有一个逆序，即 \( \iota_3 \) 位于 \( \iota_2 \) 之前。
\end{example}

\begin{example} \label{ex:ThreeInversions}
这里有三个逆序：
\begin{equation*}
  \begin{mat}[r]
    0  &0  &1  \\
    0  &1  &0  \\
    1  &0  &0
  \end{mat}
  =
  \begin{mat}
    \iota_3  \\
    \iota_2  \\
    \iota_1
  \end{mat}
\end{equation*}
\( \iota_3 \) 位于 \( \iota_1 \) 之前，
\( \iota_3 \) 位于 \( \iota_2 \) 之前，且 \( \iota_2 \) 位于
\( \iota_1 \) 之前。
\end{example}

\begin{lemma} \label{le:SwapsChangeSgn}
%<*lm:SwapsChangeSgn>
置换矩阵中的一次行对换会使逆序的个数
从偶数变为奇数，或从奇数变为偶数。
%</lm:SwapsChangeSgn>
\end{lemma}

\begin{proof}
%<*pf:SwapsChangeSgn0>
考虑对第 $j$ 行与第 $k$ 行的对换，其中 $k>j$。

若这两行相邻
\begin{equation*}
   P_{\phi}=
   \begin{mat}
     \vdots           \\
     \iota_{\phi(j)}  \\
     \iota_{\phi(k)}  \\
     \vdots
   \end{mat}
 \grstep{\rho_k\leftrightarrow\rho_j}
   \begin{mat}
     \vdots           \\
     \iota_{\phi(k)}  \\
     \iota_{\phi(j)}  \\
     \vdots
   \end{mat}
\end{equation*}
那么由于涉及这对之外其他行的逆序不受影响，
这次对换使逆序的总数改变一，
究竟是消去一个逆序还是产生一个逆序，
取决于 \( \phi(j)>\phi(k) \) 是否成立。
因此，逆序的总数从奇数变为偶数，
或从偶数变为奇数。
%</pf:SwapsChangeSgn0>

%<*pf:SwapsChangeSgn1>
若这两行不相邻，则我们可以通过一系列相邻行的对换来交换它们，
先把第~$k$ 行向上移
\begin{equation*}
   \begin{mat}
     \vdots             \\
     \iota_{\phi(j)}    \\
     \iota_{\phi(j+1)}  \\
     \iota_{\phi(j+2)}  \\
     \vdots             \\
     \iota_{\phi(k)}    \\
     \vdots
   \end{mat}
  \grstep{\rho_k\swap\rho_{k-1}}
  \repeatedgrstep{\rho_{k-1}\swap\rho_{k-2}}
  \cdots
  \grstep{\rho_{j+1}\swap\rho_j}
   \begin{mat}
     \vdots             \\
     \iota_{\phi(k)}    \\
     \iota_{\phi(j)}    \\
     \iota_{\phi(j+1)}  \\
     \vdots             \\
     \iota_{\phi(k-1)}  \\
     \vdots
   \end{mat}
\end{equation*}
%</pf:SwapsChangeSgn1>
%<*pf:SwapsChangeSgn2>
再把第~$j$ 行向下移。
\begin{equation*}
  \grstep{\rho_{j+1}\swap\rho_{j+2}}
  \repeatedgrstep{\rho_{j+2}\swap\rho_{j+3}}
  \cdots
  \grstep{\rho_{k-1}\swap\rho_k}
   \begin{mat}
     \vdots             \\
     \iota_{\phi(k)}    \\
     \iota_{\phi(j+1)}  \\
     \iota_{\phi(j+2)}  \\
     \vdots             \\
     \iota_{\phi(j)}    \\
     \vdots
   \end{mat}
\end{equation*}
这些相邻对换中的每一次都使逆序的个数
从奇数变为偶数或从偶数变为奇数。
对换的总次数 \( (k-j)+(k-j-1) \) 是奇数。
于是，总体上逆序的个数
从偶数变为奇数，或从奇数变为偶数。
%</pf:SwapsChangeSgn2>
\end{proof}

\begin{corollary}  \label{cor:ParityInversEqParitySwaps}
%<*co:ParityInversEqParitySwaps>
若置换矩阵有奇数个逆序，则把它换成
单位矩阵需要奇数次对换。
若它有偶数个逆序，则换成单位矩阵需要偶数次对换。
%</co:ParityInversEqParitySwaps>
\end{corollary}

\begin{proof}
%<*pf:ParityInversEqParitySwaps>
单位矩阵有零个逆序。
把奇数变成零需要奇数次对换，
把偶数变成零需要偶数次对换。
%</pf:ParityInversEqParitySwaps>
\end{proof}

\begin{example}
\nearbyexample{ex:ThreeInversions} 中的矩阵
用一次对换 $\rho_1\leftrightarrow\rho_3$ 就能变成单位矩阵。
（所以对换的次数不必与逆序的个数相同，
但这两个数的奇偶性是一样的。）
\end{example}

\begin{definition} \label{df:Signum}
%<*df:Signum>
置换的\definend{符号}\index{permutation!signum}\index{signum}%
\index{sgn!see {signum}}
\( \sgn(\phi) \)
在 \( \phi \) 中逆序个数为奇数时是 \( -1 \)，为偶数时是
\( +1 \)。
%</df:Signum>
\end{definition}

\begin{example}
使用 \nearbyexample{ex:AllTwoThreePerms} 中
$3$ 元置换的记号，我们有
\begin{equation*}
  P_{\phi_1}=
  \begin{mat}
    1  &0  &0  \\
    0  &1  &0  \\
    0  &0  &1
  \end{mat}
  \qquad
  P_{\phi_2}=
  \begin{mat}
    1  &0  &0  \\
    0  &0  &1  \\
    0  &1  &0  
  \end{mat}
\end{equation*}
所以
\( \sgn(\phi_1)=1 \)，因为没有逆序，
而 \( \sgn(\phi_2)=-1 \)，因为有一个逆序。
\end{example}

我们仍未证明行列式函数是良定义的，
因为除了行对换之外，
我们还没有考虑置换矩阵上的其他行变换。
我们将巧妙地绕过这个问题。
%<*DefiningDFunction>
定义函数
\( \map{d}{\matspace_{\nbyn{n}}}{\Re} \)：
修改置换展开公式，
把 $\deter{P_\phi}$ 换成 $\sgn(\phi)$。
\begin{equation*}
  d(T)=
  \sum_{\text{permutations\ }\phi}t_{1,\phi(1)}t_{2,\phi(2)}\cdots t_{n,\phi(n)}
                                 \cdot\sgn(\phi)
\end{equation*}
% This gives the same value as the permutation expansion
% because the corollary shows that $\det(P_\phi)=\sgn(\phi)$.
这个公式的优点是逆序的个数
显然是良定义的\Dash  直接数出来即可。
因此，当我们证明了~$d$
满足行列式定义中的条件时，
也就证明了一个~$\nbyn{n}$阶行列式函数的
存在性。
%</DefiningDFunction>

\begin{lemma} \label{lm:DeterminantsExist}
\index{determinant!exists}
%<*lm:DeterminantsExist>
上面的函数 \( d \) 是一个行列式。
从而 $\det_{\nbyn{n}}$ 对每个 $n$ 都存在。
%</lm:DeterminantsExist>
\end{lemma}

\begin{proof}
%<*pf:DeterminantsExist0>
我们必须验证它满足行列式定义中的四个条件，
即 \nearbydefinition{def:Det}。
%</pf:DeterminantsExist0>

%<*pf:DeterminantsExist1>
条件~(4) 容易验证：设 $I$ 是 $\nbyn{n}$ 阶单位矩阵，
在
\begin{equation*}
  d(I)=
  \sum_{\text{perm\ }\phi}
    \iota_{1,\phi(1)}\iota_{2,\phi(2)}\cdots \iota_{n,\phi(n)}
                                 \sgn(\phi)
\end{equation*}
中，求和式里所有的项都为零，
只有置换~$\phi$ 为恒等的那一项例外，
这一项给出沿对角线的乘积，
而该乘积为一。
%</pf:DeterminantsExist1>

%<*pf:DeterminantsExist2>
对于条件~(3)，设 \( T\smash[b]{\grstep{k\rho_i}}\hat{T} \)，
考虑 $d(\hat{T})$。
\begin{multline*}
  \sum_{\text{perm\ }\phi}\!\!
    \hat{t}_{1,\phi(1)}
     \cdots\hat{t}_{i,\phi(i)}\cdots\hat{t}_{n,\phi(n)}
                                 \sgn(\phi)   \\
  =\sum_{\phi}
     t_{1,\phi(1)}\cdots kt_{i,\phi(i)}\cdots t_{n,\phi(n)}
                                 \sgn(\phi)
\end{multline*}
提出因子~\( k \) 便得到所要的等式。
\begin{equation*}
  =k\cdot\sum_{\phi}
     t_{1,\phi(1)}\cdots t_{i,\phi(i)}\cdots t_{n,\phi(n)}
                                 \sgn(\phi)                 
  =k\cdot d(T)
\end{equation*}
%</pf:DeterminantsExist2>

%<*pf:DeterminantsExist3>
对于~(2)，设
\( T\smash[b]{\grstep{\rho_i\swap\rho_j}}\hat{T} \)。
我们必须证明 $d(\hat{T})$ 是 $d(T)$ 的相反数。
\begin{equation*}
  d(\hat{T})=
  \sum_{\text{perm\ }\phi}\!\!
    \hat{t}_{1,\phi(1)}
    \cdots\hat{t}_{i,\phi(i)}
    \cdots\hat{t}_{j,\phi(j)}
    \cdots \hat{t}_{n,\phi(n)}
                                 \sgn(\phi)
    \tag{$*$}
\end{equation*}
我们将证明（$*$）中的每一项都对应 $d(T)$ 中的一项，
且这两项互为相反数。
考虑 $d(\hat{T})$ 的多重线性展开式中给出项
    $\hat{t}_{1,\phi(1)}
    \cdots\hat{t}_{i,\phi(i)}
    \cdots\hat{t}_{j,\phi(j)}
    \cdots \hat{t}_{n,\phi(n)}
                                  \sgn(\phi)$ 的矩阵。
\begin{equation*}
  \begin{mat}
     \ &                  &\vdots &               &     \\
       &\hat{t}_{i,\phi(i)}  &       &               &    \\
       &                &\vdots &                 &  \\
       &                &       &\hat{t}_{j,\phi(j)} &  \\
       &                &\vdots                   &
  \end{mat}
\end{equation*}
%</pf:DeterminantsExist3>
%<*pf:DeterminantsExist4>
它是对矩阵
施加 $\rho_i\swap\rho_j$ 操作的结果。
\begin{equation*}
  \begin{mat}
       \ &              &\vdots &           &        \\
         &              &       &t_{i,\phi(j)} &  \\
         &              &\vdots &           &       \\
         &t_{j,\phi(i)}    &       &          &        \\
         &              &\vdots &          &
  \end{mat}
\end{equation*}
也就是说，带帽 $t$ 的项对应于来自 $d(T)$ 展开式的这一项：
 \(    t_{1,\sigma(1)}
       \cdots t_{j,\sigma(j)}
       \cdots t_{i,\sigma(i)}
       \cdots t_{n,\sigma(n)} \sgn(\sigma)
       \)，其中置换
\( \sigma \)
与 \( \phi \) 相同，只是把第 \( i \) 个数与第 \( j \) 个数互换：
\( \sigma(i)=\phi(j) \) 且 \( \sigma(j)=\phi(i) \)。
这两项的各因子相同：
$\hat{t}_{1,\phi(1)}=t_{1,\sigma(1)}$，\ldots，
包括来自对换过的两行的元素
$\hat{t}_{i,\phi(i)}=t_{j,\phi(i)}=t_{j,\sigma(j)}$
和 $\hat{t}_{j,\phi(j)}=t_{i,\phi(j)}=t_{i,\sigma(i)}$。 
但这两项互为相反数，
因为由 \nearbylemma{le:SwapsChangeSgn} 有 \( \sgn(\phi)=-\sgn(\sigma) \)。
%</pf:DeterminantsExist4>

%<*pf:DeterminantsExist5>
现在，任意置换 $\phi$ 都可以由某个
其他置换 $\sigma$ 通过这样的对换得到，
而且方式有且仅有一种。
因此（$*$）中的求和实际上是对所有置换的求和，
每个置换恰好取一次。
\begin{align*}
  d(\hat{T})
  &=
  \sum_{\text{perm\ }\phi}\!\!
    \hat{t}_{1,\phi(1)}
    \cdots\hat{t}_{i,\phi(i)}
    \cdots\hat{t}_{j,\phi(j)}
    \cdots \hat{t}_{n,\phi(n)}
                                 \sgn(\phi)   \\
  &=\sum_{\text{perm\ }\sigma}\!\!
     t_{1,\sigma(1)}
    \cdots t_{j,\sigma(j)}
    \cdots t_{i,\sigma(i)}
    \cdots t_{n,\sigma(n)}
                                 \cdot\bigl(-\sgn(\sigma)\bigr)        
\end{align*}
于是 \( d(\hat{T})=-d(T) \)。
%</pf:DeterminantsExist5>

%<*pf:DeterminantsExist6>
最后，对于条件~(1)，设
\( T\smash[b]{\grstep{k\rho_i+\rho_j}}\hat{T} \)。
\begin{align*}
  d(\hat{T})
  &=\sum_{\text{perm\ }\phi}
    \hat{t}_{1,\phi(1)}\cdots\hat{t}_{i,\phi(i)}
    \cdots\hat{t}_{j,\phi(j)}\cdots\hat{t}_{n,\phi(n)}
                                 \sgn(\phi)                \\
  &=\sum_{\phi}
    t_{1,\phi(1)}\cdots t_{i,\phi(i)}
    \cdots (kt_{i,\phi(j)}+t_{j,\phi(j)})\cdots t_{n,\phi(n)}
                                 \sgn(\phi)
\end{align*}
% (notice:~that's 
% \( kt_{i,\phi(j)} \), not \( kt_{j,\phi(j)} \)).
对 $kt_{i,\phi(j)}+t_{j,\phi(j)}$ 中的加法使用分配律。
\begin{equation*}
  = 
      \begin{aligned}[t]
        &\smash[b]{\sum_{\smash{\phi}}}\big[t_{1,\phi(1)}\cdots t_{i,\phi(i)}
          \cdots kt_{i,\phi(j)}\cdots t_{n,\phi(n)}
                                 \sgn(\phi)         \\ % [-1ex]
        &\hbox{}\qquad\quad\hbox{}
           +t_{1,\phi(1)}\cdots t_{i,\phi(i)}
           \cdots t_{j,\phi(j)}\cdots t_{n,\phi(n)}
                                 \sgn(\phi)\big]
      \end{aligned}                                               
\end{equation*}
把它拆成两个和式。
\begin{equation*}
  =\begin{aligned}[t] 
         &\smash[b]{\sum_{\smash{\phi}}} \:t_{1,\phi(1)}\cdots t_{i,\phi(i)}
            \cdots kt_{i,\phi(j)}\cdots t_{n,\phi(n)}
                                 \sgn(\phi)           \\ % [-1ex]               
         &\hbox{}\qquad\quad\hbox{}
            +\smash[b]{\sum_{\phi}}\:
            t_{1,\phi(1)}\cdots t_{i,\phi(i)}
            \cdots t_{j,\phi(j)}\cdots t_{n,\phi(n)}
                                 \sgn(\phi)        
     \end{aligned}                                               
\end{equation*}
识别出第二个和式。
\begin{equation*}
  =\begin{aligned}[t]
       &k\cdot \smash[b]{\sum_{\smash{\phi}}}\:
            t_{1,\phi(1)}\cdots t_{i,\phi(i)}
            \cdots t_{i,\phi(j)}\cdots t_{n,\phi(n)}
                                 \sgn(\phi)          \\  % [-.75ex]
       &\hbox{}\qquad\quad\hbox{}
             +d(T)          
     \end{aligned}
\end{equation*}
%</pf:DeterminantsExist6>
%<*pf:DeterminantsExist7>
考虑各项
\( t_{1,\phi(1)}\cdots t_{i,\phi(i)} \cdots t_{i,\phi(j)}\cdots t_{n,\phi(n)}
   \sgn(\phi)  \)。
注意下标；元素是 \( t_{i,\phi(j)} \)，而不是 \( t_{j,\phi(j)} \)。
这些项的和是矩阵~\( S \) 的行列式，
它等于 \( T \)，
只是 $S$ 的第~$j$ 行是 $T$ 的第~$i$ 行的副本，
也就是说，$S$ 有两行相同。
用与证明 \nearbylemma{le:IdenRowsDetZero} 相同的方式可以看出
$d(S)=0$：交换 $S$ 的相同的两行会改变
$d(S)$ 的符号，但由于矩阵经这次交换并不改变，
$d(S)$ 的值也必须不变，因此这个值必为零。
%</pf:DeterminantsExist7>
\end{proof}

至此我们证明了每个阶数~$\nbyn{n}$ 的行列式函数都存在。
我们已经知道，对每个阶数，行列式至多有一个。
因此，对每个阶数，行列式函数有且仅有一个。

我们以证明来自前一小节的另一个尚待证明的结论
来结束本小节。

\begin{theorem} \label{th:DetMatrixEqualsDetTrans}
\index{transpose!determinant}
%<*th:DetMatrixEqualsDetTrans>
矩阵的行列式等于其转置的行列式。
%</th:DetMatrixEqualsDetTrans>
\end{theorem}

\begin{proof}
%<*pf:DetMatrixEqualsDetTrans0>
最好通过作出一般的
$\nbyn{3}$ 情形来理解这个证明。
该论证适用于 $\nbyn{n}$ 情形，这一点将是显然的。

比较矩阵~$T$ 的置换展开
\begin{align*}
  \begin{vmat}
    t_{1,1}  &t_{1,2}  &t_{1,3}  \\
    t_{2,1}  &t_{2,2}  &t_{2,3}  \\
    t_{3,1}  &t_{3,2}  &t_{3,3}  
  \end{vmat}
  &=
  t_{1,1}t_{2,2}t_{3,3}
  \begin{vmat}
    1 &0 &0 \\
    0 &1 &0 \\
    0 &0 &1
  \end{vmat}
  +t_{1,1}t_{2,3}t_{3,2}
  \begin{vmat}
    1 &0 &0 \\
    0 &0 &1 \\
    0 &1 &0
  \end{vmat}            \\
  &\quad+
  t_{1,2}t_{2,1}t_{3,3}
  \begin{vmat}
    0 &1 &0 \\
    1 &0 &0 \\
    0 &0 &1
  \end{vmat}
  +t_{1,2}t_{2,3}t_{3,1}
  \begin{vmat}
    0 &1 &0 \\
    0 &0 &1 \\
    1 &0 &0
  \end{vmat}           \\
  &\quad+
  t_{1,3}t_{2,1}t_{3,2}
  \begin{vmat}
    0 &0 &1 \\
    1 &0 &0 \\
    0 &1 &0
  \end{vmat}
  +t_{1,3}t_{2,2}t_{3,1}
  \begin{vmat}
    0 &0 &1 \\
    0 &1 &0 \\
    1 &0 &0
  \end{vmat}
\end{align*}
与其转置的置换展开。
%</pf:DetMatrixEqualsDetTrans0>
%<*pf:DetMatrixEqualsDetTrans1>
\begin{align*}
  \begin{vmat}
    t_{1,1}  &t_{2,1}  &t_{3,1}  \\
    t_{1,2}  &t_{2,2}  &t_{3,2}  \\
    t_{1,3}  &t_{2,3}  &t_{3,3}  
  \end{vmat}
  &=
  t_{1,1}t_{2,2}t_{3,3}
  \begin{vmat}
    1 &0 &0 \\
    0 &1 &0 \\
    0 &0 &1
  \end{vmat}
  +t_{1,1}t_{3,2}t_{2,3}
  \begin{vmat}
    1 &0 &0 \\
    0 &0 &1 \\
    0 &1 &0
  \end{vmat}            \\
  &\quad+
  t_{2,1}t_{1,2}t_{3,3}
  \begin{vmat}
    0 &1 &0 \\
    1 &0 &0 \\
    0 &0 &1
  \end{vmat}
  +t_{2,1}t_{3,2}t_{1,3}
  \begin{vmat}
    0 &1 &0 \\
    0 &0 &1 \\
    1 &0 &0
  \end{vmat}           \\
  &\quad+
  t_{3,1}t_{1,2}t_{2,3}
  \begin{vmat}
    0 &0 &1 \\
    1 &0 &0 \\
    0 &1 &0
  \end{vmat}
  +t_{3,1}t_{2,2}t_{1,3}
  \begin{vmat}
    0 &0 &1 \\
    0 &1 &0 \\
    1 &0 &0
  \end{vmat}
\end{align*}
先比较六个 $t$ 的乘积。
$T$ 展开式中的那些乘积与转置展开式中的相同；
例如，$t_{1,2}t_{2,3}t_{3,1}$ 在上面，
而 $t_{3,1}t_{1,2}t_{2,3}$ 在下面。
这完全合理\Dash 上面的六个来自从~$T$ 的每一行、每一列
各取一个元素的所有方式，
而下面的六个是从~$T$ 的每一列、每一行
各取一个元素的所有方式，
所以它们当然是同一个集合。
%</pf:DetMatrixEqualsDetTrans1>

%<*pf:DetMatrixEqualsDetTrans2>
接下来观察到，在这两个展开式中，
每个 $t$ 乘积表达式所对应的
置换矩阵不一定相同。
例如，在上面，$t_{1,2}t_{2,3}t_{3,1}$ 对应于映射
$1\mapsto 2$, $2\mapsto 3$, $3\mapsto 1$ 的矩阵。
在下面，$t_{3,1}t_{1,2}t_{2,3}$
对应于映射
$1\mapsto 3$, $2\mapsto 1$, $3\mapsto 2$ 的矩阵。
第二个映射是第一个的逆映射。
这也完全合理\Dash 两者
矩阵转置与逆映射都把 $1,2$ 换成 $2,1$，
把 $2,3$ 换成 $3,2$，并把 $3,1$ 换成 $1,3$。
 
最后我们指出，$P_{\phi}$ 的行列式等于
$P_{\phi^{-1}}$ 的行列式，
正如 \nearbyexercise{exer:PermInvFacts} %
所表明的。
%</pf:DetMatrixEqualsDetTrans2>
\end{proof}


\begin{exercises}
  \item[]\textit{这里总结本书中 $2$-\hbox{} 与 $3$-置换所用的记号。
    \begin{center}
      \begin{tabular}[t]{c|cc}
        $i$          &$1$      &$2$    \\
        \hline
        $\phi_1(i)$  &$1$      &$2$     \\
        $\phi_2(i)$  &$2$      &$1$     
      \end{tabular}
      \qquad
      \begin{tabular}[t]{c|ccc}
        $i$          &$1$     &$2$   &$3$    \\
        \hline
        $\phi_1(i)$  &$1$     &$2$   &$3$    \\
        $\phi_2(i)$  &$1$     &$3$   &$2$    \\
        $\phi_3(i)$  &$2$     &$1$   &$3$    \\
        $\phi_4(i)$  &$2$     &$3$   &$1$    \\
        $\phi_5(i)$  &$3$     &$1$   &$2$    \\
        $\phi_6(i)$  &$3$     &$2$   &$1$    
      \end{tabular}
    \end{center}  }
  \item 
    写出一般 $\nbyn{2}$ 矩阵及其转置的置换展开。
    \begin{answer}
      这是 $\nbyn{2}$ 矩阵的行列式的置换展开
      \begin{equation*}
        \begin{vmat}
          a  &b  \\
          c  &d
        \end{vmat}
         =ad\cdot \begin{vmat}[r]
                    1  &0  \\
                    0  &1
                  \end{vmat}
         +bc\cdot \begin{vmat}[r]
                    0  &1 \\
                    1  &0
                  \end{vmat}
      \end{equation*}
      以及它的转置的行列式的置换展开。
      \begin{equation*}
         \begin{vmat}
            a  &c  \\
            b  &d
         \end{vmat}
         =ad\cdot\begin{vmat}[r]
                    1  &0  \\
                    0  &1
                 \end{vmat}
         +cb\cdot\begin{vmat}[r]
                   0  &1 \\
                   1  &0
                 \end{vmat}
       \end{equation*}
       与小节中描述的 $\nbyn{3}$ 展开一样，对应项的置换矩阵互为转置
       （尽管每个矩阵都是自转置的这一事实把这一点掩盖了）。
    \end{answer}
  \recommended \item 
    \textit{本题也出现在前面的小节中。}
    \begin{exparts}
      \partsitem 求每个 $2$-置换的逆。
      \partsitem 求每个 $3$-置换的逆。
    \end{exparts}
    \begin{answer}
      每一个都容易验证。
      \begin{exparts}
        \partsitem 
          \begin{tabular}[t]{r|cc}
            \textit{置换} &$\phi_1$  &$\phi_2$ \\
             \hline
            \textit{逆}     &$\phi_1$  &$\phi_2$ 
          \end{tabular}
        \partsitem 
          \begin{tabular}[t]{r|cccccc}
            \textit{置换} 
              &$\phi_1$ &$\phi_2$ &$\phi_3$ &$\phi_4$ &$\phi_5$ &$\phi_6$ \\
            \hline
            \textit{逆}
              &$\phi_1$ &$\phi_2$ &$\phi_3$ &$\phi_5$ &$\phi_4$ &$\phi_6$ 
          \end{tabular}
      \end{exparts}
    \end{answer}
  \recommended \item
    \begin{exparts}
      \partsitem 求每个 $2$-置换的符号。
      \partsitem 求每个 $3$-置换的符号。
    \end{exparts}
    \begin{answer}
      \begin{exparts}
        \partsitem \( \sgn(\phi_1)=+1 \), \( \sgn(\phi_2)=-1 \)
        \partsitem \( \sgn(\phi_1)=+1 \), \( \sgn(\phi_2)=-1 \),
              \( \sgn(\phi_3)=-1 \), \( \sgn(\phi_4)=+1 \),
              \( \sgn(\phi_5)=+1 \), \( \sgn(\phi_6)=-1 \)
      \end{exparts}  
     \end{answer}
  \item 
    找出此矩阵的置换展开中唯一非零的项。
    \begin{equation*}
      \begin{vmat}[r]
        0  &1  &0  &0  \\
        1  &0  &1  &0  \\
        0  &1  &0  &1  \\
        0  &0  &1  &0
      \end{vmat}
    \end{equation*}
    通过求相应置换的符号来计算该行列式。
    \begin{answer}
      要在置换展开中得到非零项，我们必须取 \( 1,2 \) 处的元素与
      \( 4,3 \) 处的元素。
      选定这两处之后，我们还必须取 \( 2,1 \) 处的元素与 \( 3,4 \) 处的元素。
      \( \sequence{2,1,4,3} \) 的符号是 \( +1 \)，因为由
      \begin{equation*}
        \begin{mat}[r]
          0  &1  &0  &0  \\
          1  &0  &0  &0  \\
          0  &0  &0  &1  \\
          0  &0  &1  &0            
        \end{mat}
      \end{equation*}
      经过两次行对换 $\rho_1\leftrightarrow\rho_2$ 与
      $\rho_3\leftrightarrow\rho_4$ 即可得到单位矩阵。  
    \end{answer}
  \item  
    \cite{Strang}
    $n$-置换 \( \phi=\sequence{n,n-1,\dots,2,1} \) 的符号是什么？
    \begin{answer}
      规律如下。
      \begin{center}
         \begin{tabular}[b]{c|ccccccc}
           \( i \) &\( 1 \) &\( 2 \) &\( 3 \) &\( 4 \) &\( 5 \)
              &\( 6 \) &\ldots \\
           \hline
           \( \sgn(\phi_i) \) &\( +1 \) &\( -1 \) &\( -1 \) &\( +1 \)
              &\( +1 \) &\( -1 \) &\ldots
         \end{tabular}
      \end{center}  
      于是为求 $\phi_{n!}$ 的符号，我们用 $n!-1$ 减一，再看它除以四的余数。
      若余数是 $1$ 或 $2$，则符号为 $-1$；否则为 $+1$。
      对 $n>4$，数 $n!$ 能被四整除，所以 $n!-1$ 除以四的余数是
      $-1$（更准确地说，余数是 $3$），因此符号是 $+1$。
      $n=1$ 的情形符号为 $+1$，$n=2$ 的情形符号为 $-1$，
      $n=3$ 的情形符号为 $-1$。
    \end{answer}
  \item \label{exer:PermInvFacts}
     证明这些结论。   
     \begin{exparts}
        \partsitem 每个置换都有逆。
        \partsitem \( \sgn(\phi^{-1})=\sgn(\phi) \)
        \partsitem 每个置换都是另一个置换的逆。
     \end{exparts}
    \begin{answer}
      \begin{exparts}
       \partsitem 我们可以把置换看成
         \( \map{\phi}{\set{1,\ldots,n}}{\set{1,\ldots,n}} \)
         这样单射且满射的映射。
         任何单射且满射的映射都有逆。
       \partsitem 若从 \( P_\phi \)
         变到单位矩阵总需要奇数次对换，那么从单位矩阵变到
         \( P_\phi \) 也总需要奇数次对换（任何对换都是可逆的）。
       \partsitem 这就是第一问。
      \end{exparts}  
    \end{answer}
  \item \label{exer:PermInvIffMatTrans}
      证明对任意置换 $\phi$，置换逆的矩阵就是置换的矩阵的转置
      $P_{\phi^{-1}}=\trans{P_{\phi}}$。
      \begin{answer}
        若 $\phi(i)=j$ 则 $\phi^{-1}(j)=i$。
        于是注意到 $P_{\phi}$ 在 $i,j$ 处有一个 $1$ 当且仅当 $\phi(i)=j$，
        而 $P_{\phi^{-1}}$ 在 $j,i$ 处有一个 $1$ 当且仅当 $\phi^{-1}(j)=i$，
        结论即得。
      \end{answer}
  \recommended \item 
    证明有 \( m \) 个逆序的置换矩阵可以经过 \( m \) 步行对换变成单位矩阵。
    将此与\nearbycorollary{cor:ParityInversEqParitySwaps}对比。
    \begin{answer}
      这并不是说 \( m \) 是变成单位矩阵所需的最少对换次数，
      也不是说 \( m \) 是最多的。
      它说的是存在一种恰好经 \( m \) 步对换变成单位矩阵的方法。

      设 \( \iota_j \) 是相对于前面某行来说第一个出现逆序的行，
      并设 \( \iota_k \) 是给出该逆序的第一行。
      我们有这样一个行区间。
      \begin{equation*}
         \begin{mat}
           \vdots      \\
           \iota_k     \\
           \iota_{r_1} \\
           \vdots      \\
           \iota_{r_s} \\
           \iota_j     \\
           \vdots
         \end{mat}
         \qquad j<k<r_1<\dots <r_s
      \end{equation*}
      作对换。
      \begin{equation*}
         \begin{mat}
           \vdots      \\
           \iota_j     \\
           \iota_{r_1} \\
           \vdots      \\
           \iota_{r_s} \\
           \iota_k     \\
           \vdots
         \end{mat}
      \end{equation*}
      第二个矩阵的逆序少了一个，因为该区间内的逆序少了一个
      （\( s \) 对 \( s+1 \)），而涉及区间之外各行的逆序不受影响。

      如此进行下去，每作一次行对换就把逆序的个数减少一个。
      当逆序不复存在时，结果就是单位矩阵。

      与\nearbycorollary{cor:ParityInversEqParitySwaps}的对比在于，
      本练习的陈述是一个“存在”命题：
      存在一种恰好经~$m$ 步对换变成单位矩阵的方法。
      而该推论是一个“对所有”命题：对变成单位矩阵的所有对换方式，
      奇偶性（是偶还是奇）都相同。 
    \end{answer}
  \recommended \item 
    对任意置换 $\phi$，令 $g(\phi)$ 是如下定义的整数。
    \begin{equation*}
       g(\phi)=\prod_{i<j} [\phi(j)-\phi(i)]
    \end{equation*}
    （这是在所有满足 \( i<j \) 的指标 \( i \) 与 \( j \) 上，
    对上述形式的项作的乘积。）
    \begin{exparts}
      \partsitem 计算 $g$ 在所有 $2$-置换上的值。
      \partsitem 计算 $g$ 在所有 $3$-置换上的值。
      \partsitem 证明 $g(\phi)$ 不为 $0$。
      \partsitem 证明这个结论。
         \begin{equation*}
            \sgn(\phi)=\frac{g(\phi)}{\absval{g(\phi)}}
         \end{equation*}
    \end{exparts}
    许多作者把这个公式当作符号函数的定义。
    \begin{answer}
      \begin{exparts}
        \partsitem 首先，$g(\phi_1)$ 是单独一个因子 $2-1$ 的乘积，
          所以 $g(\phi_1)=1$。
          其次，$g(\phi_2)$ 是单独一个因子 $1-2$ 的乘积，
          所以 $g(\phi_2)=-1$。
        \partsitem 
            \begin{tabular}[t]{r|cccccc}
               \textit{置换 $\phi$} 
                 &$\phi_{1}$ &$\phi_{2}$ &$\phi_{3}$ 
                    &$\phi_{4}$ &$\phi_{5}$ &$\phi_{6}$ \\
               \hline
               $g(\phi)$ &$2$ &$-2$ &$-2$ &$2$ &$2$ &$-2$ 
            \end{tabular}
        \partsitem 它是一些非零项的乘积。
        \partsitem 注意 \( \phi(j)-\phi(i) \) 为负当且仅当
           \( \iota_{\phi(j)} \) 与 \( \iota_{\phi(i)} \) 处于它们通常顺序的一个逆序之中。
      \end{exparts}  
     \end{answer}
%  \item Prove that the signum function on permutations is unique in that
%    if \( f \) is a map from the permutations of \( \set{1,\ldots,n} \)
%    to the integers that is multiplicative:
%    \( f(\sigma\phi)=f(\sigma)f(\phi) \) then \( f \) is either the
%    function that always gives zero, or the function that always gives
%    one, or is the signum.
%    \begin{answer}
%      Consider the identity permutation.
%      \begin{equation*}
%        \map{\identity}{\set{1,2,\ldots,n}}{\set{1,2,\ldots,n}}    
%      \end{equation*}
%      Note that $f(\identity)$ equals $0$ or $1$ because
%      the multiplicative property gives 
%      $f(\identity)
%       =f(\composed{\identity}{\identity})
%       =f(\identity)\cdot f(\identity)
%       =f(\identity)^2$.
%      If $f(\identity)=0$ then $f$ is the constant function 
%      $f(\phi)=0$ because 
%      $f(\phi)=f(\composed{\identity}{\phi})=0\cdot f(\phi)=0$.
%
%      If $f(\identity)=1$ then consider the permutation
%      $\tau$ that just swaps $i$ and $j$
%      \begin{equation*}
%        1\mapsto 1
%        \quad 2\mapsto 2
%        \quad \cdots 
%        \quad i\mapsto j
%        \quad \cdots
%        \quad j\mapsto i
%        \quad \cdots
%        \quad n\mapsto n
%      \end{equation*}
%      (the $n=1$ special case is easy).
%      Note that $\composed{\tau}{\tau}=\identity$ so 
%      $1=f(\composed{\tau}{\tau})=f(\tau)\cdot f(\tau)$,
%      and $f(\tau)$ is either $+1$ or $-1$.
%    \end{answer}
\end{exercises}



\index{determinant|)}
